\subsection{Hall correspondences: product, coproduct, primitives}
\label{subsec:hall-correspondences}

The third entry of the Dirac--Igusa pro-object is the Hall structure
on the cosection-reduced d-critical heart.  It is the algebraic
content that the orientation lane equips and the hybrid carrier
indexes.  Three operations live here: the Hall product
$m_R^{\mathrm{red}}$, the collision coproduct $\Delta_R^{\mathrm{ch}}$,
and the primitive projection $P_R^{\Prim}$.  After normal-ordered
Gram descent these form at most a graded bialgebra-and-primitive
datum.  The word ``Hopf'' is reserved for the additional antipode data
of Definition~\ref{def:finite-hall-antipode-datum}; recognition with
the Borcherds--Kac--Moody target is the \S5.6 obligation only after
that data and the radical quotient data have been supplied.

The Kontsevich--Soibelman cohomological Hall algebra of a
$3$-Calabi--Yau category \cite{KSCoHA} is the prototype; on
$K3\times E$ the cosection reduction \cite{JoyceSong,JoyceUpmeier}
replaces the unreduced critical CoHA, and the wrapped sector is the
residual the hybrid carrier of \S\ref{subsec:hybrid-ran-carrier}
retains.  The Maulik--Okounkov stable-envelope construction and the
Schiffmann--Vasserot theorems on Drinfeld doubles of CoHA operate
on the same input; their integral form is the Hall side of the
present construction.

\begin{definition}[Finite compact Hall object]
\label{def:finite-geometric-hall-object}
Fix a finite HN height \(R\) and a finite retained charge set
\(\widehat\Gamma_R\).  A finite compact Hall object at height \(R\) is
the \(\widehat\Gamma_R\)-graded object
\[
\mathcal H_R^{\mathrm{geom}}
=
\bigoplus_{\widehat c\in\widehat\Gamma_R}
H_*^{\mathrm{BM}}\!\left(
\mathfrak M^{\mathrm{red}}_{R,\widehat c},
\Phi^{\mathrm{red}}_{R,\widehat c}\otimes
o_{R,\widehat c}\right),
\]
where the summand is defined only after the following data are fixed:
the finite-type retained rigidified stack
\(\mathfrak M^{\mathrm{red}}_{R,\widehat c}\) with finite residual
inertia, the cosection-reduced BBDJS vanishing-cycle complex
\(\Phi^{\mathrm{red}}_{R,\widehat c}\), the Joyce--Upmeier orientation
line \(o_{R,\widehat c}\), and an admissible compact-support
Borel--Moore six-functor model giving finite-rank groups.  The parity
is the cohomological parity of these Borel--Moore groups together with
the orientation parity recorded in the source basis provenance.

The object \(\mathcal H_R^{\mathrm{geom}}\) is only the finite
coefficient object.  It does not include the Hall product, coproduct,
unit, counit, Hopf pairing, radical quotient, PBW filtration, primitive
subspace, normal-ordered Gram descent, or transition maps.  Those are
additional source rows.  A scalar trace, Euler characteristic, signed
root multiplicity, Pfaffian product, target root window, or denominator
product is not a substitute for the displayed Borel--Moore direct sum.
\end{definition}

The packet
\(\texttt{certificates/hall/finite\_geometric\_hall\_object}\)
verifies
\(\texttt{FINITE\_GEOMETRIC\_HALL\_OBJECT\_DEFINITION\_VERIFIED}\):
it records Definition~\ref{def:finite-geometric-hall-object} as the row
\(341\) definition and imports the retained finite-moduli packet, the
reduced-orientation obstruction ledger, the transition
vanishing-cycle obstruction ledger, and the compact Hall source
obstruction ledger.  The present source packets do not yet populate
\(\mathcal H_R^{\mathrm{geom}}\): no orientation-line rows,
vanishing-cycle coefficient rows, object-component rows, source-degree
rows, homogeneous basis provenance rows, or finite-rank rows are
supplied.  Hence row \(341\) is a definition of the finite compact Hall
object, not a proof that the compact \(K3\times E\) source stage has
been populated.

\subsubsection*{The compact Hall product}

\begin{definition}[Compact Hall product on the cosection-reduced
heart]
\label{def:compact-hall-product}
Let $\mathcal C_{\sigma,S,\le R}$ be the finite HN quotient of
\S\ref{subsec:hybrid-ran-carrier}.  Assume the retained Liu--Hilbert
compactified correspondence and the compact-support pushforward of
Proposition~\ref{prop:finite-dcritical-hall-product}, including finite
residual inertia, base-change, projection-formula, Thom--Sebastiani, and
quotient-after-correspondence coherences.  The compact Hall
product at height $R$ is the operation
\[
m_R^{\mathrm{red}}=
q_{c,c'}{}_!\circ\operatorname{TS}^{\mathrm{red}}_{c,c'}\circ p_{c,c'}^*:
H_*^{\mathrm{BM}}\!\bigl(\mathfrak M_c^\sigma\times\mathfrak M_{c'}^\sigma,
\Phi_c^{\mathrm{red}}\boxtimes\Phi_{c'}^{\mathrm{red}}
\otimes o_c\boxtimes o_{c'}\bigr)
\to
H_*^{\mathrm{BM}}\!\bigl(\mathfrak M_{c+c'}^\sigma,
\Phi_{c+c'}^{\mathrm{red}}\otimes o_{c+c'}\bigr)
\]
attached to the retained compactified extension correspondence
\[
\mathfrak M_c^\sigma\times\mathfrak M_{c'}^\sigma
\xleftarrow{p_{c,c'}}\overline{\mathfrak E}_{c,c'}^{\mathrm{ret},\sigma}
\xrightarrow{q_{c,c'}}\mathfrak M_{c+c'}^\sigma,
\]
in the cosection-reduced d-critical heart, with vanishing-cycle
coefficients $\Phi^{\mathrm{red}}$, Joyce--Upmeier orientation lines
$o$, and Thom--Sebastiani transport
$\operatorname{TS}^{\mathrm{red}}_{c,c'}$.  This definition names the
operator once the three functorial pieces \(p^*\),
\(\operatorname{TS}^{\mathrm{red}}\), and \(q_!\) are supplied in the
chosen six-functor model.  It does not prove that \(p^*\) exists, that
\(q_!\) exists for the retained correspondence, or that \(q_!\) is
independent of a compactification; these are separate functorial
obligations.
It also does not supply product matrices, associativity, transition
compatibility, or primitive closure.  Associativity on the finite stage
is the eight-word two-step binary flag pentagon together with the
quotient-after-correspondence and HN-transition coherences of
\S\ref{subsec:hybrid-ran-carrier}.
\end{definition}

The packet
\(\texttt{certificates/hall/finite\_hall\_product\_operator}\)
verifies
\(\texttt{FINITE\_HALL\_PRODUCT\_OPERATOR\_DEFINITION\_VERIFIED}\):
it records the row \(342\) definition
\[
m_R^{\mathrm{red}}=q_!\operatorname{TS}^{\mathrm{red}}p^*
\]
and imports the row \(341\) finite compact Hall object, the
compact-support exceptional pushforward model, the compact Hall source
obstruction ledger, and the transition-Hall-product obstruction ledger.
The present packets give only the formal product expression and the
chosen compact-support notation.  They do not populate retained
extension-correspondence rows, \(p^*\) rows, \(q_!\) rows, reduced
Thom--Sebastiani rows on the Hall correspondence, product matrix rows,
base-change or projection-formula rows, compactification-independence
rows, associativity rows, or product-transition rows.

\begin{proposition}[Pullback leg of the finite Hall product]
\label{prop:finite-hall-product-pullback-existence}
\textnormal{[six-functor criterion proved; retained \(p\)-row not supplied]}
Let
\[
p_{c,c'}:\overline{\mathfrak E}_{R,c,c'}^{\mathrm{ret}}
\longrightarrow
\mathfrak M_{R,c}^{\mathrm{red}}\times
\mathfrak M_{R,c'}^{\mathrm{red}}
\]
be a supplied retained extension-correspondence morphism at finite HN
height \(R\).  Assume that \(p_{c,c'}\) is a finite-type morphism in
the retained Artin-stack category on which the chosen constructible
six-functor formalism is installed, and that the input coefficient
\[
K_{R,c,c'}=
(\Phi^{\mathrm{red}}_{R,c}\otimes o_{R,c})
\boxtimes
(\Phi^{\mathrm{red}}_{R,c'}\otimes o_{R,c'})
\]
is a bounded constructible complex on the input product stack.  Then
the pullback
\[
p_{c,c'}^*K_{R,c,c'}
\in
D^b_c(\overline{\mathfrak E}_{R,c,c'}^{\mathrm{ret}})
\]
exists in the same six-functor formalism.  This is the \(p^*\)-leg of
the finite Hall product \(m_R^{\mathrm{red}}\).

The present packets do not supply the hypotheses for the actual
compact \(K3\times E\) Hall product.  The finite-moduli packet has no
extension-stack rows, finite-type flag rows, source-map rows, or
d-critical coefficient rows on the extension correspondence.  The row
\(341\) finite-Hall-object packet has no object-component or
coefficient rows.  The row \(342\) product-operator packet has no
populated \(p^*\) rows.  The compact-support packet gives notation for
compact-support pushforwards, not a populated pullback leg.  Hence row
\(343\) is presently this pullback-existence criterion and obstruction
ledger, not a proof that the retained Hall product has a populated
\(p^*\)-leg.
\end{proposition}

\begin{proof}
In a constructible six-functor theory on the retained finite-type
Artin-stack category, every morphism \(p\) in the category carries a
pullback functor
\[
p^*:D^b_c(Y)\longrightarrow D^b_c(X).
\]
The finite-type and constructibility hypotheses ensure that bounded
constructible coefficients remain in the retained coefficient category
after pullback.  Applying this to the supplied extension morphism
\(p_{c,c'}\) and to \(K_{R,c,c'}\) gives the displayed object.  The
argument uses only the pullback functor; it does not construct
\(q_!\), prove Thom--Sebastiani compatibility, produce product
matrices, or prove associativity.
\end{proof}

The packet
\(\texttt{certificates/hall/finite\_hall\_product\_pullback\_functor}\)
verifies
\(\texttt{FINITE\_HALL\_PRODUCT\_PULLBACK\_FUNCTOR\_OBSTRUCTION\_VERIFIED}\):
it imports the row \(342\) product-operator packet, the finite-moduli
obstruction ledger, the row \(341\) finite-Hall-object packet, and the
compact-support model; it records the pullback-existence criterion; and
it keeps the retained extension correspondence, finite-type \(p\)-map,
constructible input coefficient, populated \(p^*\)-row, and product
matrix rows open.

\begin{proposition}[Exceptional pushforward leg of the finite Hall
product]
\label{prop:finite-hall-product-pushforward-existence}
\textnormal{[compact-support criterion proved; retained \(q\)-row not
supplied]}
Let
\[
q_{c,c'}:\overline{\mathfrak E}_{R,c,c'}^{\mathrm{ret}}
\longrightarrow
\mathfrak M_{R,c+c'}^{\mathrm{red}}
\]
be a supplied retained extension-correspondence morphism at finite HN
height \(R\), and let
\[
K_{\mathfrak E}\in
D^b_c(\overline{\mathfrak E}_{R,c,c'}^{\mathrm{ret}})
\]
be the bounded constructible coefficient obtained after the preceding
pullback and reduced Thom--Sebastiani transport.  Assume that
\(q_{c,c'}\) lies in the chosen retained finite-type Artin-stack
category with finite residual inertia and that one of the following
two pieces of functorial data is supplied:
\begin{enumerate}
\item \(q_{c,c'}\) is proper in that category; or
\item a compact-support model
\[
j_q:\overline{\mathfrak E}_{R,c,c'}^{\mathrm{ret}}
\hookrightarrow
\overline{\mathfrak E}_{q},\qquad
\overline q:\overline{\mathfrak E}_{q}\to
\mathfrak M_{R,c+c'}^{\mathrm{red}},
\qquad
q_{c,c'}=\overline q\circ j_q
\]
is supplied as in
Definition~\ref{def:compact-support-exceptional-pushforward-model},
with \(j_{q,!}K_{\mathfrak E}\) and
\(\overline q_*j_{q,!}K_{\mathfrak E}\) in the same constructible
six-functor formalism.
\end{enumerate}
Then the exceptional pushforward leg
\[
q^{\mathrm{cs}}_{c,c',!}K_{\mathfrak E}
:=
\overline q_*j_{q,!}K_{\mathfrak E}
\in D^b_c(\mathfrak M_{R,c+c'}^{\mathrm{red}})
\]
exists; in the proper case it is the ordinary proper pushforward
\(q_{c,c',*}K_{\mathfrak E}\).  This is the \(q_!\)-leg of the finite
Hall product \(m_R^{\mathrm{red}}\).

The present packets do not supply the hypotheses for the actual
compact \(K3\times E\) Hall product.  The row \(342\) product operator
records \(q_!\) only as notation.  The row \(343\) packet proves only
the pullback criterion.  The finite-moduli ledger has no extension
flag-stack row, no retained \(q\)-map row, no properness row for
\(q\), and no d-critical coefficient row on the extension
correspondence.  The compact-support model defines
\(\overline q_*j_{q,!}\), but its map-population and
choice-independence tables are empty.  The compact Hall source packet
has no product matrix rows.  Hence row \(344\) is this exceptional
pushforward criterion and obstruction ledger, not a proof that the
retained Hall product has a populated \(q_!\)-leg.
\end{proposition}

\begin{proof}
In the proper case, the six-functor formalism supplies the proper
direct image
\[
q_{c,c',*}:D^b_c(\overline{\mathfrak E}_{R,c,c'}^{\mathrm{ret}})
\longrightarrow
D^b_c(\mathfrak M_{R,c+c'}^{\mathrm{red}}),
\]
and properness preserves constructibility.  In the compact-support
case, the supplied open immersion gives extension by zero
\[
j_{q,!}K_{\mathfrak E}\in D^b_c(\overline{\mathfrak E}_{q}),
\]
and the supplied proper map \(\overline q\) gives the proper direct
image
\[
\overline q_*j_{q,!}K_{\mathfrak E}
\in D^b_c(\mathfrak M_{R,c+c'}^{\mathrm{red}}).
\]
This composite is exactly the compact-support exceptional pushforward
of Definition~\ref{def:compact-support-exceptional-pushforward-model}.
The argument does not compare two compactifications, prove base
change or the projection formula, construct the reduced
Thom--Sebastiani row, produce product matrices, or prove associativity.
\end{proof}

The packet
\(\texttt{certificates/hall/finite\_hall\_product\_pushforward\_functor}\)
verifies
\(\texttt{FINITE\_HALL\_PRODUCT\_PUSHFORWARD\_FUNCTOR\_OBSTRUCTION\_VERIFIED}\):
it imports the row \(342\) product-operator packet, the row \(343\)
pullback packet, the finite-moduli obstruction ledger, the row \(341\)
finite-Hall-object packet, the compact-support model, and the compact
Hall source obstruction ledger; it records the compact-support
pushforward criterion; and it keeps the retained extension
correspondence, retained \(q\)-map, properness or compactification
row, constructible coefficient row, populated \(q_!\)-row,
compactification independence, and product matrix rows open.

\begin{proposition}[Compactification independence criterion for the
Hall \(q_!\)-leg]
\label{prop:finite-hall-product-qbang-compactification-independence}
\textnormal{[common-refinement criterion proved; compactification pair
not supplied]}
Let
\[
q:\mathfrak E\longrightarrow Z
\]
be a supplied retained Hall target morphism at finite HN height \(R\),
and let \(K_{\mathfrak E}\in D^b_c(\mathfrak E)\) be the bounded
constructible coefficient on the correspondence source.  Suppose two
compact-support models for \((q,K_{\mathfrak E})\) are supplied:
\[
j_i:\mathfrak E\hookrightarrow\overline{\mathfrak E}_i,\qquad
\overline q_i:\overline{\mathfrak E}_i\to Z,\qquad
q=\overline q_i\circ j_i,\qquad i=1,2,
\]
with \(\overline q_i\) proper.  Assume further that a common proper
refinement is supplied:
\[
j_{12}:\mathfrak E\hookrightarrow\overline{\mathfrak E}_{12},
\qquad
r_i:\overline{\mathfrak E}_{12}\to\overline{\mathfrak E}_i,\qquad
\overline q_{12}:\overline{\mathfrak E}_{12}\to Z,
\]
where the maps \(r_i\) and \(\overline q_{12}\) are proper and satisfy
\[
r_i\circ j_{12}=j_i,\qquad
\overline q_i\circ r_i=\overline q_{12}.
\]
Finally assume that the extension-by-zero comparison morphisms
\[
r_{i,*}j_{12,!}K_{\mathfrak E}
\xrightarrow{\;\sim\;}
j_{i,!}K_{\mathfrak E},
\qquad i=1,2,
\]
are isomorphisms in the chosen constructible six-functor formalism.
Then the two compact-support pushforwards are canonically identified:
\[
\overline q_{1,*}j_{1,!}K_{\mathfrak E}
\cong
\overline q_{12,*}j_{12,!}K_{\mathfrak E}
\cong
\overline q_{2,*}j_{2,!}K_{\mathfrak E}.
\]
Thus \(q_!^{\mathrm{cs}}K_{\mathfrak E}\) is independent of the two
compactifications whenever the common-refinement datum and the two
comparison isomorphisms are supplied.

The present packets do not supply the hypotheses for the actual
compact \(K3\times E\) Hall product.  The compact-support model has
only the abstract formula
\(\overline q_*j_{q,!}\); its choice-independence table is empty.  The
row \(344\) packet records the existence criterion for a supplied
compactification, not a comparison between two choices.  The
finite-moduli ledger has no retained \(q\)-map row, no compactification
pair row, and no common-refinement row.  Hence row \(345\) is this
common-refinement criterion and obstruction ledger, not a proof that
the actual retained Hall \(q_!\)-leg is already independent of
compactification.
\end{proposition}

\begin{proof}
For \(i=1,2\), proper functoriality gives
\[
\overline q_{i,*}r_{i,*}
=
(\overline q_i\circ r_i)_*
=
\overline q_{12,*}.
\]
Composing this equality with the supplied comparison isomorphism
\[
r_{i,*}j_{12,!}K_{\mathfrak E}\simeq j_{i,!}K_{\mathfrak E}
\]
gives
\[
\overline q_{i,*}j_{i,!}K_{\mathfrak E}
\simeq
\overline q_{i,*}r_{i,*}j_{12,!}K_{\mathfrak E}
=
\overline q_{12,*}j_{12,!}K_{\mathfrak E}.
\]
The two choices \(i=1,2\) are therefore identified through the same
middle object.  The proof uses only proper functoriality and the
supplied extension-by-zero comparison maps; it does not construct the
common refinement, prove base change or the projection formula, produce
product matrices, or prove associativity.
\end{proof}

The packet
\(\texttt{certificates/hall/finite\_hall\_product\_compactification\_independence}\)
verifies
\(\texttt{FINITE\_HALL\_PRODUCT\_COMPACTIFICATION\_INDEPENDENCE\_OBSTRUCTION\_VERIFIED}\):
it imports the row \(342\) product-operator packet, the row \(344\)
pushforward packet, the finite-moduli obstruction ledger, the
compact-support model, and the compact Hall source obstruction ledger;
it records the common-refinement criterion; and it keeps the retained
\(q\)-map, compactification-pair row, common-refinement row,
comparison-isomorphism row, populated choice-independence row, and
product matrix rows open.

\begin{proposition}[Associativity criterion via retained two-step
flags]
\label{prop:finite-hall-product-associativity-two-step-flags}
\textnormal{[two-step-flag criterion proved; populated \(m_R\)-rows
not supplied]}
Let \(c_1,c_2,c_3\in\widehat\Gamma_R\).  Suppose the finite Hall
product legs
\[
m_{a,b}^{\mathrm{red}}
=q_{a,b,!}^{\mathrm{cs}}\operatorname{TS}^{\mathrm{red}}_{a,b}p_{a,b}^*
\]
are supplied, choice-independent, and constructible for the four pairs
\[
(c_1,c_2),\quad (c_1+c_2,c_3),\quad
(c_2,c_3),\quad (c_1,c_2+c_3)
\]
that occur in the two parenthesizations.  Assume also that a retained
two-step extension flag stack
\[
\mathfrak F^{\mathrm{ret}}_{R;c_1,c_2,c_3}
\]
is supplied, with comparison morphisms to the two iterated
correspondence fibre products
\[
\lambda:\mathfrak F^{\mathrm{ret}}_{R;c_1,c_2,c_3}\to
\overline{\mathfrak E}^{\mathrm{ret}}_{R;c_1+c_2,c_3}
\times_{\mathfrak M^{\mathrm{red}}_{R,c_1+c_2}}
\overline{\mathfrak E}^{\mathrm{ret}}_{R;c_1,c_2},
\]
\[
\rho:\mathfrak F^{\mathrm{ret}}_{R;c_1,c_2,c_3}\to
\overline{\mathfrak E}^{\mathrm{ret}}_{R;c_1,c_2+c_3}
\times_{\mathfrak M^{\mathrm{red}}_{R,c_2+c_3}}
\overline{\mathfrak E}^{\mathrm{ret}}_{R;c_2,c_3}.
\]
Assume the squares defining \(\lambda\) and \(\rho\) satisfy the
base-change isomorphisms for the relevant \(p^*\)- and \(q_!^{cs}\)
legs, that the projection formula holds for the coefficient objects,
and that the reduced Thom--Sebastiani identifications on the two
parenthesizations agree on
\(\mathfrak F^{\mathrm{ret}}_{R;c_1,c_2,c_3}\).  Finally assume that
the two compact-support target pushforwards from the flag stack to
\(\mathfrak M^{\mathrm{red}}_{R,c_1+c_2+c_3}\) are identified by the
compactification-independence criterion of
Proposition~\ref{prop:finite-hall-product-qbang-compactification-independence}.

Then for all supplied constructible inputs
\[
K_i\in D^b_c(\mathfrak M^{\mathrm{red}}_{R,c_i}),\qquad i=1,2,3,
\]
there is a canonical associativity isomorphism
\[
m^{\mathrm{red}}_{c_1+c_2,c_3}
\bigl(m^{\mathrm{red}}_{c_1,c_2}(K_1,K_2),K_3\bigr)
\cong
m^{\mathrm{red}}_{c_1,c_2+c_3}
\bigl(K_1,m^{\mathrm{red}}_{c_2,c_3}(K_2,K_3)\bigr).
\]
After Borel--Moore homology and after choosing homogeneous bases, this
is the rowwise associativity identity for the finite product matrix
\(M_R\).

The present packets do not supply the hypotheses for the actual
compact \(K3\times E\) Hall product.  Rows \(342\)--\(345\) define the
operator and record criteria for \(p^*\), \(q_!\), and
compactification independence, but their populated product,
pushforward, and independence tables are empty.  The hybrid two-step
flag packet constructs the eight local/wrapped flag stacks and
comparison maps before the reduced quotient; the wordwise
associativity packets are conditional on functorial rows and do not
populate the compact Hall carrier.  The compact Hall source packet has
no product matrix entries and no Hall-bialgebra associativity rows.
Hence row \(346\) is this associativity criterion and obstruction
ledger, not a proof that the actual retained \(m_R\) is already
associative as a populated finite matrix.
\end{proposition}

\begin{proof}
Expand the left parenthesization by the definition of
\(m^{\mathrm{red}}\).  It is the compact-support pushforward along the
left iterated correspondence after pulling back
\(K_1\boxtimes K_2\boxtimes K_3\) and applying reduced
Thom--Sebastiani twice.  The right parenthesization is the analogous
pushforward along the right iterated correspondence.  Pull both
iterated correspondences back to
\(\mathfrak F^{\mathrm{ret}}_{R;c_1,c_2,c_3}\) by \(\lambda\) and
\(\rho\).  The supplied base-change and projection-formula
isomorphisms identify the two iterated pull--push expressions with
compact-support pushforwards from this common flag stack.  The supplied
Thom--Sebastiani associativity isomorphism identifies the two
coefficient complexes on the flag stack.  The supplied
compactification-independence comparison identifies the two final
compact-support pushforwards to
\(\mathfrak M^{\mathrm{red}}_{R,c_1+c_2+c_3}\).  The resulting
isomorphism is the displayed associativity isomorphism.  Passing to
Borel--Moore homology and to homogeneous bases gives the matrix
identity whenever the product matrix rows are supplied.
\end{proof}

The packet
\(\texttt{certificates/hall/finite\_hall\_product\_associativity\_two\_step\_flags}\)
verifies
\(\texttt{FINITE\_HALL\_PRODUCT\_ASSOCIATIVITY\_TWO\_STEP\_FLAGS\_OBSTRUCTION\_VERIFIED}\):
it imports the row \(342\) product operator, rows \(344\)--\(345\), the
finite-moduli obstruction ledger, the compact-support model, the
eight-word two-step flag-stack packet, the conditional wordwise
associativity packets, and the compact Hall source obstruction ledger;
it records the two-step-flag associativity criterion; and it keeps the
populated compact product rows, compact Hall two-step flag rows,
functorial base-change/projection/Thom--Sebastiani rows, quotient
descent, transition rows, product matrices, and Hall-bialgebra
associativity rows open.

For local-only inputs and ordered mixed and wrapped inputs, the same
construction produces the three operation families
$\mu_R^{\mathrm{loc}}$, $\mu_R^{\mathrm{hyb}}$, $\mu_R^{\mathrm{wr}}$
of \S\ref{subsec:hybrid-ran-carrier}.  Existence of the product as a
genuine morphism on Borel--Moore chains requires the
admissibility/properness residuals of that section to vanish.

\subsubsection*{The collision coproduct on the source coalgebra}

The collision coproduct is defined on the same correspondence
geometry, read in the opposite direction: instead of contracting an
extension to a single middle term, one expands an object into the pairs of
its sub/quotient extensions.  At finite stage, this is the
$2$-step flag stack already supplied as part of the hybrid datum.

\begin{definition}[Collision coproduct]
\label{def:collision-coproduct}
The finite collision coproduct
\[
\Delta_R^{\mathrm{ch}}:C_{X,R}\longrightarrow C_{X,R}\otimes C_{X,R}
\]
is the chain-level pull-push along the $2$-step flag stack
$\mathfrak{Flag}^{(2)}_{c,c',R}$ of subobjects of charge $c$ and
quotients of charge $c'$ in $\mathfrak M_{c+c',R}^\sigma$, computed
in the cosection-reduced d-critical heart, before the reduced
$E$-quotient.  If the higher-coloured residual, unit maps, refinement
coherences, and quotient-after-correspondence residuals of
\S\ref{subsec:hybrid-ran-carrier} vanish, then after $Q_{E,R}$ it
satisfies the coalgebra identities
\[
(\Delta_R^{\mathrm{ch}}\otimes 1)\Delta_R^{\mathrm{ch}}
=(1\otimes\Delta_R^{\mathrm{ch}})\Delta_R^{\mathrm{ch}},
\qquad
(\varepsilon_R^C\otimes 1)\Delta_R^{\mathrm{ch}}
=(1\otimes\varepsilon_R^C)\Delta_R^{\mathrm{ch}}=\mathrm{id},
\]
which are the dual of the Hall pentagon on the same flag stack.  For
$R'\ge R$ the transition map $\rho^C_{R'R}$ commutes with
$\Delta^{\mathrm{ch}}$.
\end{definition}

\begin{definition}[Compact Hall coproduct on the cosection-reduced
heart]
\label{def:compact-hall-coproduct}
Let
\[
\mathfrak M^{\mathrm{red}}_{R,c}\times
\mathfrak M^{\mathrm{red}}_{R,c'}
\xleftarrow{\ p_{c,c'}\ }
\overline{\mathfrak E}^{\mathrm{ret}}_{R,c,c'}
\xrightarrow{\ q_{c,c'}\ }
\mathfrak M^{\mathrm{red}}_{R,c+c'}
\]
be the retained compactified extension correspondence at finite HN
height \(R\), with the same reduced vanishing-cycle and orientation
coefficients as in Definition~\ref{def:compact-hall-product}.  The
compact Hall coproduct operator is the formal pull--inverse
Thom--Sebastiani--push operator
\[
\Delta_R^{\mathrm{red}}
=
p_{c,c',!}^{\mathrm{cs}}\circ
(\operatorname{TS}^{\mathrm{red}}_{c,c'})^{-1}\circ
q_{c,c'}^*.
\]
It sends a coefficient on
\(\mathfrak M^{\mathrm{red}}_{R,c+c'}\) to a coefficient on
\(\mathfrak M^{\mathrm{red}}_{R,c}\times
\mathfrak M^{\mathrm{red}}_{R,c'}\), and after Borel--Moore
Kunneth identification it is read as
\[
\Delta_R^{\mathrm{red}}:
\mathcal H^{\mathrm{geom}}_{R,c+c'}
\longrightarrow
\mathcal H^{\mathrm{geom}}_{R,c}\otimes
\mathcal H^{\mathrm{geom}}_{R,c'}.
\]
This definition names the operator once the three functorial pieces
\(q^*\), \((\operatorname{TS}^{\mathrm{red}})^{-1}\), and
\(p_!^{\mathrm{cs}}\) are supplied in the chosen six-functor model.  It
does not prove that \(q^*\) exists, that \(p_!\) exists for the
reversed correspondence, that \(p_!\) is independent of a
compactification, or that the coproduct is coassociative.  It also does
not supply counit maps, coproduct matrices, transition compatibility,
bialgebra compatibility, or primitive closure.
\end{definition}

The packet
\(\texttt{certificates/hall/finite\_hall\_coproduct\_operator}\)
verifies
\(\texttt{FINITE\_HALL\_COPRODUCT\_OPERATOR\_DEFINITION\_VERIFIED}\):
it records the row \(347\) definition
\[
\Delta_R^{\mathrm{red}}
=p_!(\operatorname{TS}^{\mathrm{red}})^{-1}q^*
\]
and imports the row \(341\) finite compact Hall object, the
compact-support exceptional pushforward model, the compact Hall source
obstruction ledger, and the transition-Hall-coproduct obstruction
ledger.  The present packets give only the formal coproduct expression
and the chosen compact-support notation.  They do not populate retained
splitting-correspondence rows, \(q^*\) rows, \(p_!\) rows, inverse
Thom--Sebastiani rows, compactification-independence rows, coproduct
matrix rows, counit rows, coassociativity rows, bialgebra rows, or
coproduct-transition rows.

\begin{proposition}[Exceptional pushforward leg of the finite Hall
coproduct]
\label{prop:finite-hall-coproduct-pushforward-existence}
\textnormal{[compact-support criterion proved; retained \(p\)-row not
supplied]}
Let
\[
p_{c,c'}:\overline{\mathfrak E}_{R,c,c'}^{\mathrm{ret}}
\longrightarrow
\mathfrak M_{R,c}^{\mathrm{red}}\times
\mathfrak M_{R,c'}^{\mathrm{red}}
\]
be a supplied retained splitting-correspondence morphism at finite HN
height \(R\), and let
\[
K_{\mathfrak E}\in
D^b_c(\overline{\mathfrak E}_{R,c,c'}^{\mathrm{ret}})
\]
be the bounded constructible coefficient obtained after the preceding
pullback \(q_{c,c'}^*\) and inverse reduced Thom--Sebastiani
transport.  Assume that \(p_{c,c'}\) lies in the chosen retained
finite-type Artin-stack category with finite residual inertia and that
one of the following two pieces of functorial data is supplied:
\begin{enumerate}
\item \(p_{c,c'}\) is proper in that category; or
\item a compact-support model
\[
j_p:\overline{\mathfrak E}_{R,c,c'}^{\mathrm{ret}}
\hookrightarrow
\overline{\mathfrak E}_{p},\qquad
\overline p:\overline{\mathfrak E}_{p}\to
\mathfrak M_{R,c}^{\mathrm{red}}\times
\mathfrak M_{R,c'}^{\mathrm{red}},
\qquad
p_{c,c'}=\overline p\circ j_p
\]
is supplied as in
Definition~\ref{def:compact-support-exceptional-pushforward-model},
with \(j_{p,!}K_{\mathfrak E}\) and
\(\overline p_*j_{p,!}K_{\mathfrak E}\) in the same constructible
six-functor formalism.
\end{enumerate}
Then the exceptional pushforward leg
\[
p^{\mathrm{cs}}_{c,c',!}K_{\mathfrak E}
:=
\overline p_*j_{p,!}K_{\mathfrak E}
\in
D^b_c(\mathfrak M_{R,c}^{\mathrm{red}}\times
\mathfrak M_{R,c'}^{\mathrm{red}})
\]
exists; in the proper case it is the ordinary proper pushforward
\(p_{c,c',*}K_{\mathfrak E}\).  This is the \(p_!\)-leg of the finite
Hall coproduct \(\Delta_R^{\mathrm{red}}\).

The present packets do not supply the hypotheses for the actual
compact \(K3\times E\) Hall coproduct.  The row \(347\) coproduct
operator records \(p_!\) only as notation.  The finite-moduli ledger
has no extension flag-stack row, no retained splitting \(p\)-map row,
no properness row for \(p\), and no d-critical coefficient row on the
splitting correspondence.  The compact-support model defines
\(\overline p_*j_{p,!}\), but its map-population and
choice-independence tables are empty.  The compact Hall source packet
has no coproduct matrix rows, and the transition-Hall-coproduct packet
has no coproduct-transport rows.  Hence row \(348\) is this
exceptional pushforward criterion and obstruction ledger, not a proof
that the retained Hall coproduct has a populated \(p_!\)-leg.
\end{proposition}

\begin{proof}
In the proper case, the six-functor formalism supplies the proper
direct image
\[
p_{c,c',*}:
D^b_c(\overline{\mathfrak E}_{R,c,c'}^{\mathrm{ret}})
\longrightarrow
D^b_c(\mathfrak M_{R,c}^{\mathrm{red}}\times
\mathfrak M_{R,c'}^{\mathrm{red}}),
\]
and properness preserves constructibility.  In the compact-support
case, the supplied open immersion gives extension by zero
\[
j_{p,!}K_{\mathfrak E}\in D^b_c(\overline{\mathfrak E}_{p}),
\]
and the supplied proper map \(\overline p\) gives the proper direct
image
\[
\overline p_*j_{p,!}K_{\mathfrak E}
\in
D^b_c(\mathfrak M_{R,c}^{\mathrm{red}}\times
\mathfrak M_{R,c'}^{\mathrm{red}}).
\]
This composite is exactly the compact-support exceptional pushforward
of Definition~\ref{def:compact-support-exceptional-pushforward-model}.
The argument does not compare two compactifications, prove base
change or the projection formula, construct the \(q^*\) row, construct
the inverse reduced Thom--Sebastiani row, produce coproduct matrices,
or prove coassociativity.
\end{proof}

The packet
\(\texttt{certificates/hall/finite\_hall\_coproduct\_pushforward\_functor}\)
verifies
\(\texttt{FINITE\_HALL\_COPRODUCT\_PUSHFORWARD\_FUNCTOR\_OBSTRUCTION\_VERIFIED}\):
it imports the row \(347\) coproduct-operator packet, the
finite-moduli obstruction ledger, the row \(341\) finite-Hall-object
packet, the compact-support model, the compact Hall source obstruction
ledger, and the transition-Hall-coproduct obstruction ledger; it
records the compact-support pushforward criterion; and it keeps the
retained splitting correspondence, retained \(p\)-map, properness or
compactification row, constructible coefficient row, inverse
Thom--Sebastiani row, populated \(p_!\)-row, compactification
independence, coproduct matrix rows, and transition-coproduct rows
open.

\begin{proposition}[Coassociativity criterion via retained two-step
splitting flags]
\label{prop:finite-hall-coproduct-coassociativity-two-step-flags}
\textnormal{[two-step-splitting-flag criterion proved; populated
\(\Delta_R\)-rows not supplied]}
Let \(c_1,c_2,c_3\in\widehat\Gamma_R\).  Suppose the finite Hall
coproduct legs
\[
\Delta_{a,b}^{\mathrm{red}}
=p_{a,b,!}^{\mathrm{cs}}
(\operatorname{TS}_{a,b}^{\mathrm{red}})^{-1}q_{a,b}^*
\]
are supplied, choice-independent, and constructible for the four
pairs
\[
(c_1+c_2,c_3),\quad (c_1,c_2),\quad
(c_1,c_2+c_3),\quad (c_2,c_3)
\]
that occur in the two coproduct parenthesizations.  Assume also that
a retained two-step splitting flag stack
\[
\mathfrak F^{\mathrm{ret}}_{R;c_1,c_2,c_3}
\]
is supplied, with comparison morphisms to the two iterated splitting
correspondence fibre products
\[
\lambda:\mathfrak F^{\mathrm{ret}}_{R;c_1,c_2,c_3}\to
\overline{\mathfrak E}^{\mathrm{ret}}_{R;c_1+c_2,c_3}
\times_{\mathfrak M^{\mathrm{red}}_{R,c_1+c_2}}
\overline{\mathfrak E}^{\mathrm{ret}}_{R;c_1,c_2},
\]
\[
\rho:\mathfrak F^{\mathrm{ret}}_{R;c_1,c_2,c_3}\to
\overline{\mathfrak E}^{\mathrm{ret}}_{R;c_1,c_2+c_3}
\times_{\mathfrak M^{\mathrm{red}}_{R,c_2+c_3}}
\overline{\mathfrak E}^{\mathrm{ret}}_{R;c_2,c_3}.
\]
Assume the squares defining \(\lambda\) and \(\rho\) satisfy the
base-change isomorphisms for the relevant \(q^*\)- and
\(p_!^{\mathrm{cs}}\)-legs, that the projection formula holds for the
coefficient objects, and that the inverse reduced
Thom--Sebastiani identifications on the two parenthesizations agree on
\(\mathfrak F^{\mathrm{ret}}_{R;c_1,c_2,c_3}\).  Finally assume that
the two compact-support pushforwards from the flag stack to
\[
\mathfrak M^{\mathrm{red}}_{R,c_1}\times
\mathfrak M^{\mathrm{red}}_{R,c_2}\times
\mathfrak M^{\mathrm{red}}_{R,c_3}
\]
are identified by a supplied compactification-independence comparison
for the \(p_!^{\mathrm{cs}}\)-legs.

Then for every supplied constructible input
\[
K\in D^b_c(\mathfrak M^{\mathrm{red}}_{R,c_1+c_2+c_3})
\]
there is a canonical coassociativity isomorphism
\[
(\Delta_{c_1,c_2}^{\mathrm{red}}\otimes 1)
\Delta_{c_1+c_2,c_3}^{\mathrm{red}}(K)
\cong
(1\otimes\Delta_{c_2,c_3}^{\mathrm{red}})
\Delta_{c_1,c_2+c_3}^{\mathrm{red}}(K).
\]
After Borel--Moore homology and after choosing homogeneous bases, this
is the rowwise coassociativity identity for the finite coproduct
matrix \(D_R\).

The present packets do not supply the hypotheses for the actual
compact \(K3\times E\) Hall coproduct.  Row \(347\) names the
coproduct operator but has no populated \(\Delta_R\)-rows.  Row
\(348\) proves only the \(p_!\)-existence criterion for a supplied
retained splitting map, and its pushforward table is empty.  The
finite-moduli ledger has no retained two-step splitting flag stack.
The compact-support model has no base-change, projection-formula,
Thom--Sebastiani, or choice-independence rows.  The hybrid two-step
flag packet constructs local/wrapped flag stacks before the reduced
quotient and does not populate the compact Hall carrier.  The compact
Hall source packet has no coproduct matrix entries and no
Hall-bialgebra coassociativity rows.  Hence row \(349\) is this
coassociativity criterion and obstruction ledger, not a proof that the
actual retained \(\Delta_R\) is already coassociative as a populated
finite matrix.
\end{proposition}

\begin{proof}
Expand the left parenthesization by the definition of
\(\Delta^{\mathrm{red}}\).  It is the compact-support pushforward to
\(\mathfrak M^{\mathrm{red}}_{R,c_1}\times
\mathfrak M^{\mathrm{red}}_{R,c_2}\times
\mathfrak M^{\mathrm{red}}_{R,c_3}\) obtained by pulling \(K\) back
along the two relevant \(q\)-maps, applying inverse reduced
Thom--Sebastiani twice, and pushing along the two relevant
\(p\)-maps.  The right parenthesization is the analogous expression
for the other iterated splitting correspondence.  Pull both iterated
correspondences back to
\(\mathfrak F^{\mathrm{ret}}_{R;c_1,c_2,c_3}\) by \(\lambda\) and
\(\rho\).  The supplied base-change and projection-formula
isomorphisms identify the two iterated pull--inverse-TS--push
expressions with compact-support pushforwards from this common flag
stack.  The supplied inverse Thom--Sebastiani coassociativity
isomorphism identifies the two coefficient complexes on the flag
stack.  The supplied compactification-independence comparison
identifies the two compact-support pushforwards to the triple product.
The resulting isomorphism is the displayed coassociativity
isomorphism.  Passing to Borel--Moore homology and to homogeneous
bases gives the matrix identity whenever the coproduct matrix rows are
supplied.
\end{proof}

The packet
\(\texttt{certificates/hall/finite\_hall\_coproduct\_coassociativity\_two\_step\_flags}\)
verifies
\(\texttt{FINITE\_HALL\_COPRODUCT\_COASSOCIATIVITY\_TWO\_STEP\_FLAGS\_OBSTRUCTION\_VERIFIED}\):
it imports the row \(347\) coproduct operator, row \(348\), the
finite-moduli obstruction ledger, the compact-support model, the
eight-word two-step flag-stack packet, the compact Hall source
obstruction ledger, and the transition-Hall-coproduct obstruction
ledger; it records the two-step-splitting-flag coassociativity
criterion; and it keeps the populated compact coproduct rows, compact
Hall two-step splitting flag rows, functorial
base-change/projection/inverse-Thom--Sebastiani rows,
compactification-independence rows, quotient descent, transition rows,
coproduct matrices, and Hall-bialgebra coassociativity rows open.

\begin{proposition}[Bialgebra compatibility criterion for the finite
Hall product and coproduct]
\label{prop:finite-hall-bialgebra-compatibility}
\textnormal{[product--coproduct square criterion proved; populated
bialgebra rows not supplied]}
Let \(a,b,u,v\in\widehat\Gamma_R\) with \(u+v=a+b\), and set
\[
\mathsf S(a,b;u,v)
=\{(a_1,a_2,b_1,b_2)\mid
a=a_1+a_2,\ b=b_1+b_2,\ u=a_1+b_1,\ v=a_2+b_2\}.
\]
The set \(\mathsf S(a,b;u,v)\) is finite at finite HN height.  Suppose
the product legs \(m^{\mathrm{red}}_{a,b}\),
\(m^{\mathrm{red}}_{a_1,b_1}\),
\(m^{\mathrm{red}}_{a_2,b_2}\) and the coproduct legs
\(\Delta^{\mathrm{red}}_{u,v}\),
\(\Delta^{\mathrm{red}}_{a_1,a_2}\),
\(\Delta^{\mathrm{red}}_{b_1,b_2}\) are supplied, constructible, and
choice-independent for every
\((a_1,a_2,b_1,b_2)\in\mathsf S(a,b;u,v)\).  Assume also that a
retained product--splitting square stack
\[
\mathfrak Q^{\mathrm{ret};u,v}_{R;a,b}
\]
is supplied, with comparison morphisms to the iterated correspondence
for \(\Delta_{u,v}m_{a,b}\) and to the finite disjoint union, indexed
by \(\mathsf S(a,b;u,v)\), of the iterated correspondences for
\[
(m_{a_1,b_1}^{\mathrm{red}}\otimes
m_{a_2,b_2}^{\mathrm{red}})
(1\otimes\tau\otimes1)
(\Delta_{a_1,a_2}^{\mathrm{red}}\otimes
\Delta_{b_1,b_2}^{\mathrm{red}}).
\]
Assume the comparison squares satisfy proper base change for the
relevant \(q^*\), \(p^*\), \(q_!^{\mathrm{cs}}\), and
\(p_!^{\mathrm{cs}}\) legs; the projection formula for the coefficient
objects; the product--coproduct Thom--Sebastiani coherence identifying
the product-side and coproduct-side coefficient complexes on
\(\mathfrak Q^{\mathrm{ret};u,v}_{R;a,b}\); the Koszul sign
convention for the middle braiding \(\tau\); and compact-support
choice-independence for all final pushforwards to
\(\mathfrak M^{\mathrm{red}}_{R,u}\times
\mathfrak M^{\mathrm{red}}_{R,v}\).

Then for all supplied constructible inputs
\[
K_a\in D^b_c(\mathfrak M^{\mathrm{red}}_{R,a}),\qquad
K_b\in D^b_c(\mathfrak M^{\mathrm{red}}_{R,b}),
\]
there is a canonical bialgebra-compatibility isomorphism
\[
\Delta^{\mathrm{red}}_{u,v}
m^{\mathrm{red}}_{a,b}(K_a,K_b)
\cong
\bigoplus_{\mathsf S(a,b;u,v)}
(m^{\mathrm{red}}_{a_1,b_1}\otimes
m^{\mathrm{red}}_{a_2,b_2})
(1\otimes\tau\otimes1)
(\Delta^{\mathrm{red}}_{a_1,a_2}K_a\otimes
\Delta^{\mathrm{red}}_{b_1,b_2}K_b).
\]
After Borel--Moore homology and after choosing homogeneous bases, this
is the rowwise matrix identity
\[
D_R^{a+b;u,v}M_R^{a,b}
=
\sum_{\mathsf S(a,b;u,v)}
(M_R^{a_1,b_1}\otimes M_R^{a_2,b_2})
(1\otimes\tau\otimes1)
(D_R^{a;a_1,a_2}\otimes D_R^{b;b_1,b_2}).
\]

The present packets do not supply the hypotheses for the actual
compact \(K3\times E\) Hall bialgebra.  Row \(342\) names the product
operator and row \(347\) names the coproduct operator; neither packet
has populated operation rows.  Rows \(346\) and \(349\) record the
associativity and coassociativity criteria, but their identity tables
are empty.  The finite-moduli ledger has no retained
product--splitting square stack.  The compact-support model has no
base-change, projection-formula, Thom--Sebastiani, or
choice-independence rows for the mixed product--coproduct square.  The
compact Hall source packet has no \(M\)-entries, \(D\)-entries, or
Hall-bialgebra identity rows.  Hence row \(350\) is this bialgebra
compatibility criterion and obstruction ledger, not a proof that the
actual retained finite source already satisfies
\(\Delta m=(m\otimes m)\Delta\).
\end{proposition}

\begin{proof}
Expand the left side by the definitions of \(m^{\mathrm{red}}\) and
\(\Delta^{\mathrm{red}}\).  It is a compact-support pull--TS--push
along the product correspondence followed by a
pull--inverse-TS--push along the splitting correspondence.  Expand the
right side by first splitting the two inputs, applying the Koszul
braiding to exchange the two middle factors, and multiplying the two
output pairs.  Pull the two iterated correspondences to
\(\mathfrak Q^{\mathrm{ret};u,v}_{R;a,b}\).  Proper base change and
the projection formula identify the two iterated pull-push functors on
this common square stack.  The supplied product--coproduct
Thom--Sebastiani coherence identifies the coefficient complexes; the
middle exchange contributes exactly the Koszul sign \(\tau\).  The
compact-support choice-independence comparisons identify the final
pushforwards to
\(\mathfrak M^{\mathrm{red}}_{R,u}\times
\mathfrak M^{\mathrm{red}}_{R,v}\).  The resulting isomorphism is the
displayed bialgebra-compatibility isomorphism.  Applying
Borel--Moore homology, Kunneth, and homogeneous compact-source bases
gives the displayed finite matrix identity whenever \(M_R\), \(D_R\),
and the bialgebra identity rows are supplied.
\end{proof}

The packet
\(\texttt{certificates/hall/finite\_hall\_bialgebra\_compatibility}\)
verifies
\(\texttt{FINITE\_HALL\_BIALGEBRA\_COMPATIBILITY\_OBSTRUCTION\_VERIFIED}\):
it imports the product and coproduct operator packets, the row \(346\)
associativity packet, the row \(349\) coassociativity packet, the
finite-moduli obstruction ledger, the compact-support model, the
compact Hall source obstruction ledger, and the product- and
coproduct-transition obstruction ledgers; it records the
product--coproduct square criterion; and it keeps the populated
product rows, populated coproduct rows, retained product--splitting
square rows, base-change/projection/Thom--Sebastiani rows,
compactification-independence rows, \(M\)- and \(D\)-matrices,
transition rows, and Hall-bialgebra compatibility rows open.

\begin{definition}[Finite Hall antipode datum]
\label{def:finite-hall-antipode-datum}
A finite Hall antipode datum at HN height \(R\) is extra structure on a
supplied finite Hall bialgebra
\[
(\mathcal H_R,m_R,\Delta_R,\eta_R,\epsilon_R)
\]
whose unit, counit, associativity, coassociativity, and
product--coproduct compatibility rows have zero defect.  It consists
of either:
\begin{enumerate}
\item explicit homogeneous antipode matrices
\[
S_{R,\gamma}:\mathcal H_{R,\gamma}\longrightarrow
\mathcal H_{R,\iota(\gamma)}
\]
for a supplied charge involution \(\iota\) preserving the retained
finite window; or
\item a supplied connected conilpotent filtration for which the
standard recursive antipode
\[
S_R(1)=1,\qquad
S_R(x)=-x-\sum S_R(x')\,x''
\quad
\bigl(\widetilde\Delta_Rx=\sum x'\otimes x''\bigr)
\]
terminates in every retained homogeneous block.
\end{enumerate}
In either presentation the convolution identities
\[
m_R(S_R\otimes\mathrm{id})\Delta_R=\eta_R\epsilon_R,\qquad
m_R(\mathrm{id}\otimes S_R)\Delta_R=\eta_R\epsilon_R
\]
must be supplied as zero-defect rows, and \(S_R\) must commute with the
HN transition maps on every block where the pro-object is formed.  A
finite Hall bialgebra becomes a finite Hall Hopf algebra only after
this antipode datum is supplied.

The present compact \(K3\times E\) packets do not supply an antipode
datum.  The source packet has no antipode matrix rows, no charge
involution row, no connected conilpotent recursion row, and no left or
right convolution identity row.  The bialgebra-compatibility packet
records only the criterion for
\(\Delta m=(m\otimes m)(1\otimes\tau\otimes1)(\Delta\otimes\Delta)\).
Hopf-pairing rows, Drinfeld double comparisons, bar-coalgebra
deconcatenation, primitive closure, scalar traces, and signed
multiplicities are not antipode data.
\end{definition}

The packet
\(\texttt{certificates/hall/finite\_hall\_antipode\_datum}\)
verifies
\(\texttt{FINITE\_HALL\_ANTIPODE\_DATUM\_OBSTRUCTION\_VERIFIED}\):
it imports the row \(350\) bialgebra-compatibility packet, the compact
Hall source obstruction ledger, and the product- and
coproduct-transition obstruction ledgers; it records
Definition~\ref{def:finite-hall-antipode-datum}; and it keeps the
antipode matrices, charge involution or connected-conilpotent
filtration, unit and counit rows, bialgebra rows, convolution-identity
rows, and transition-antipode rows open.  Thus row \(351\) defines the
antipode data required by any finite Hall Hopf-algebra claim; it does
not assert that the current compact Hall source is a Hopf algebra.

\begin{proposition}[No Hopf algebra without an antipode; \textit{proved}]
\label{prop:finite-hall-nonhopf-boundary}
Let
\[
(\mathcal H_R,m_R,\Delta_R,\eta_R,\epsilon_R)
\]
be a supplied finite Hall bialgebra stage.  If the finite Hall
antipode datum of Definition~\ref{def:finite-hall-antipode-datum} is
not supplied, then this stage is not a finite Hall Hopf algebra in the
verified source package of the manuscript.  It may be used only as a
bialgebra-level input.  Every downstream occurrence of a Hopf algebra
claim is therefore conditional on the antipode matrices, or on the
terminating connected-conilpotent recursion, together with the two
convolution identities and transition compatibility.

This is not a non-existence theorem for antipodes on the underlying
graded vector space.  It is the source-data boundary used in this
manuscript: absent \(S_R\) and its convolution identities, the Hopf
algebra name is unavailable.
\end{proposition}

\begin{proof}
A Hopf algebra is a bialgebra equipped with an antipode.  In the
finite Hall presentation this means exactly the datum
\(S_R\) and the two identities
\[
m_R(S_R\otimes\mathrm{id})\Delta_R=\eta_R\epsilon_R,\qquad
m_R(\mathrm{id}\otimes S_R)\Delta_R=\eta_R\epsilon_R .
\]
If no such row is supplied, no convolution inverse for the identity
endomorphism of \(\mathcal H_R\) has been constructed.  Bialgebra
compatibility, Hopf-pairing adjointness, a target Drinfeld double, a
bar-coalgebra model, primitive closure, and scalar trace data do not
define that convolution inverse.  The verified source package
therefore remains bialgebra-level until the antipode datum is added.
\end{proof}

The packet
\(\texttt{certificates/hall/finite\_hall\_nonhopf\_boundary}\)
verifies \(\texttt{FINITE\_HALL\_NONHOPF\_BOUNDARY\_VERIFIED}\): it
imports the row \(351\) antipode-datum packet, the row \(350\)
bialgebra-compatibility packet, and the compact Hall source
obstruction ledger; it records
Proposition~\ref{prop:finite-hall-nonhopf-boundary}; it checks that
the current source has no antipode rows, no convolution identity rows,
and no finite Hall Hopf-algebra claim; and it records explicitly that
no non-existence theorem for future antipodes is being asserted.

This is the source coalgebra of \S\ref{subsec:koszul-source}; its
identification with the bar construction
$\bar B_E^{\mathrm{ch}}(\mathcal F_X^{\mathrm{hyb}})$ is the bar
comparison there.

\subsubsection*{Primitive projection and the bialgebra datum}

The primitive subspace is a bialgebra notion.  It uses the
augmentation, unit, counit, and coproduct; its closure under the
supercommutator uses bialgebra compatibility.  It does not make the
source object Hopf.  Hopf-level recognition is available only after
the antipode datum named in
Proposition~\ref{prop:finite-hall-nonhopf-boundary} is supplied.

\begin{definition}[Finite primitive Hall subspace]
\label{def:finite-hall-primitive-subspace}
Let
\((\mathcal H_R,\Delta_R,\eta_R,\epsilon_R)\) be a finite counital
Hall coalgebra stage with supplied homogeneous finite bases.  Put
\[
I_R=\ker\epsilon_R .
\]
The reduced coproduct is the map
\[
\bar\Delta_R:I_R\longrightarrow I_R\otimes I_R,\qquad
\bar\Delta_R(x)=
\Delta_Rx-\eta_R(1)\otimes x-x\otimes\eta_R(1),
\]
equivalently the projection of \(\Delta_R|_{I_R}\) to
\(I_R\otimes I_R\).  The finite primitive Hall subspace is
\[
P_R=\Prim(\mathcal H_R)
=\ker\bigl(\bar\Delta_R:I_R\to I_R\otimes I_R\bigr).
\]
In finite coordinates, row \(353\) consists of the counit matrix,
the augmentation-ideal basis, the reduced-coproduct matrix
\(\bar D_R\), and kernel rows giving a basis, inclusion, and projection
for \(P_R\).  It is a definition of a subspace of the supplied finite
coalgebra.  It does not assert that the compact \(K3\times E\) source
has supplied those rows, and it does not assert primitive closure under
the Hall commutator.
\end{definition}

The packet
\(\texttt{certificates/hall/finite\_hall\_primitive\_subspace\_definition}\)
verifies
\(\texttt{FINITE\_HALL\_PRIMITIVE\_SUBSPACE\_DEFINITION\_OBSTRUCTION\_VERIFIED}\):
it imports the row \(352\) non-Hopf boundary, the row \(350\)
bialgebra-compatibility packet, and the compact Hall source
obstruction ledger; it records
Definition~\ref{def:finite-hall-primitive-subspace}; and it keeps the
source coproduct matrix, unit and counit rows, augmentation-ideal
basis, reduced-coproduct matrix, primitive-kernel rows, primitive
basis, inclusion, projection, and degree--parity split open.  Thus
row \(353\) defines \(P_R=\ker\bar\Delta_R\) for supplied finite data;
it does not populate \(P_R\) in the present compact source.

\begin{definition}[Primitive Hall sub-bracket]
\label{def:primitive-hall-bracket}
Write
\[
\widetilde P_X=H^0\Prim_{\mathrm{prot}}(\mathcal F_X)
=\ker\bigl(\widetilde\Delta:
\mathcal F_X\to\mathcal F_X\otimes\mathcal F_X\bigr)
\]
for the primitive layer of the protected Hall bialgebra object, where
$\widetilde\Delta=\Delta-(\eta\varepsilon\otimes\mathrm{id}+
\mathrm{id}\otimes\eta\varepsilon)$.  The Hall product restricts to
a $\widehat\Gamma_X$-graded bracket
\[
[\,\cdot\,,\,\cdot\,]_X:\widetilde P_{X,\widehat c}\otimes
\widetilde P_{X,\widehat c'}
\longrightarrow\widetilde P_{X,\widehat c\star\widehat c'},
\]
the protected primitive Hall bracket.  After normal-ordered Gram
descent (\S\ref{subsec:hybrid-ran-carrier}) and the Hopf-radical
quotient, the descended bracket is $\Gamma_{\mathrm{gram}}$-graded
and the descended primitive layer is the candidate $P_X^{\Pi,\mathrm{red}}$
to be identified with $\mathfrak g_{\Delta_5}$.
\end{definition}

The finite packet
\(\texttt{certificates/charge/hall\_pairing\_pushforward\_compatibility}\)
supplies only the grading calculation behind this sentence: a supplied
\(\widehat\Gamma_X\)-graded bracket and a supplied homogeneous
degree-zero pairing remain graded after pushforward by the additive
normal-ordered map \(\overline\Pi_X\).  The compact Hall
correspondence, primitive closure, Hopf adjointness, Frobenius cyclic
identity, quotient non-degeneracy, and protected trace remain separate
data.

\begin{proposition}[Product-only Hall data do not define primitives;
\textit{proved}]
\label{prop:hall-product-only-no-primitives}
Let \(\mathcal H_R\) be a finite HN stage carrying the Hall product
\(m_R\) and unit \(\eta_R\).  Without a counital coproduct
\(\Delta_R\) and the bialgebra compatibility of \(m_R\) with
\(\Delta_R\), there is no defined primitive subspace
\[
\Prim(\mathcal H_R)=
\ker\bigl(\Delta_R-\eta_R\varepsilon_R\otimes\mathrm{id}
-\mathrm{id}\otimes\eta_R\varepsilon_R\bigr),
\]
and the supercommutator of \(m_R\) does not define a Lie bracket on
primitives.  Thus a compact Hall product, even if associative and
oriented, is not a primitive-recognition datum.
\end{proposition}

\begin{proof}
The displayed formula is the definition of primitives in a counital
bialgebra.  If \(\Delta_R\) or \(\varepsilon_R\) is absent, the kernel
is not an object.  If \(\Delta_R\) is present but not compatible with
\(m_R\), the bialgebra identity
\[
\Delta_R([x,y])=[\Delta_Rx,\Delta_Ry]
\]
in the graded tensor-product algebra does not hold.  Hence
\(\Delta_R x=x\otimes1+1\otimes x\) and
\(\Delta_R y=y\otimes1+1\otimes y\) no longer imply
\(\Delta_R([x,y])=[x,y]\otimes1+1\otimes[x,y]\).

The product does not determine the coproduct.  Over a characteristic
zero field, the graded algebra \(A=k[x,y]\), \(|y|=2|x|\), has the
primitive coproduct
\(\Delta_0x=x\otimes1+1\otimes x\),
\(\Delta_0y=y\otimes1+1\otimes y\), and also the counital
coassociative algebra coproduct
\[
\Delta_1x=x\otimes1+1\otimes x,\qquad
\Delta_1y=y\otimes1+1\otimes y+x\otimes x.
\]
The multiplication is the same in both bialgebras, but
\(\Prim(A,\Delta_0)\) and \(\Prim(A,\Delta_1)\) are different
subspaces: \(y\) is primitive for \(\Delta_0\), while
\(y-\frac12x^2\) is primitive for \(\Delta_1\).  Hence the Hall
product alone cannot determine \(P_R^{\Prim}\), the Hopf radical, or
the PBW comparison used in primitive recognition.
\end{proof}

\begin{definition}[Finite Hall bialgebra defect tuple]
\label{def:finite-hall-bialgebra-defect-tuple}
Fix homogeneous finite bases on a retained HN stage and write
\((M,D,\eta,\epsilon)\) for the matrices of the Hall product,
coproduct, unit, and counit.  The finite Hall bialgebra defect tuple is
\[
\mathfrak O_{\mathrm{Hbialg},R}
=\bigl(
\mathfrak o^{\mathrm{unit},L}_R,\mathfrak o^{\mathrm{unit},R}_R,
\mathfrak o^{\mathrm{counit},L}_R,\mathfrak o^{\mathrm{counit},R}_R,
\mathfrak o^{\mathrm{assoc}}_R,\mathfrak o^{\mathrm{coassoc}}_R,
\mathfrak o^{\mathrm{bialg}}_R,\mathfrak o^{\mathrm{prim}}_R
\bigr),
\]
where
\[
\mathfrak o^{\mathrm{assoc}}_R
=m_R(m_R\otimes1)-m_R(1\otimes m_R),\qquad
\mathfrak o^{\mathrm{coassoc}}_R
=(\Delta_R\otimes1)\Delta_R-(1\otimes\Delta_R)\Delta_R,
\]
\[
\mathfrak o^{\mathrm{bialg}}_R
=\Delta_Rm_R-(m_R\otimes m_R)(1\otimes\tau\otimes1)
(\Delta_R\otimes\Delta_R),
\]
with \(\tau\) the Koszul sign braiding, and
\[
\mathfrak o^{\mathrm{prim}}_R(x,y)
=\widetilde\Delta_R
\bigl(m_R(x,y)-(-1)^{|x||y|}m_R(y,x)\bigr)
\]
for \(x,y\in\ker\widetilde\Delta_R\cap\ker\epsilon_R\).  The unit and
counit components are the four usual left/right defects
\[
m_R(\eta_R\otimes1)-1,\quad m_R(1\otimes\eta_R)-1,\quad
(\epsilon_R\otimes1)\Delta_R-1,\quad
(1\otimes\epsilon_R)\Delta_R-1.
\]
\end{definition}

\begin{proposition}[Hall primitives are formed only after the bialgebra
defects vanish; \textit{proved}]
\label{prop:hall-bialgebra-defects-primitive-closure}
At a finite HN stage, the matrices
\((M,D,\eta,\epsilon)\) define a counital graded bialgebra whose
supercommutator restricts to a bracket on the primitive subspace if and
only if the first seven components of
\(\mathfrak O_{\mathrm{Hbialg},R}\) vanish; then
\(\mathfrak o^{\mathrm{prim}}_R=0\) follows.  Conversely, a finite
source proof table in which any component of
\(\mathfrak O_{\mathrm{Hbialg},R}\) is absent or non-zero does not
define the primitive Hall Lie superalgebra used in
Theorem~\ref{thm:primitive-recognition}.
\end{proposition}

\begin{proof}
The unit, counit, associativity, coassociativity, and bialgebra
compatibility equations are exactly the first seven displayed defects.
When they vanish, the standard calculation in a graded bialgebra gives,
for primitive \(x,y\),
\[
\Delta_R([x,y])
=[x,y]\otimes1+1\otimes[x,y],
\qquad
\epsilon_R([x,y])=0,
\]
because \(\Delta_R\) is multiplicative and \(\epsilon_R\) is an algebra
homomorphism.  Hence the supercommutator lies in the primitive subspace,
which is the vanishing of \(\mathfrak o^{\mathrm{prim}}_R\).  If one of
the first seven equations is missing or non-zero, the object is not a
counital bialgebra on the chosen finite stage, and the displayed
primitive-bracket calculation is not available.
\end{proof}

The packet
\(\texttt{certificates/hall/finite\_hall\_primitive\_commutator\_closure}\)
verifies
\(\texttt{FINITE\_HALL\_PRIMITIVE\_COMMUTATOR\_CLOSURE\_OBSTRUCTION\_VERIFIED}\):
it imports the row \(353\) primitive-subspace definition, the row
\(350\) bialgebra-compatibility packet, and the compact Hall source
obstruction ledger; it records
Proposition~\ref{prop:hall-bialgebra-defects-primitive-closure} as the
row \(354\) relative proof; and it keeps the product matrix, coproduct
matrix, unit and counit rows, primitive-kernel basis, bialgebra
identity rows, primitive-closure defect rows, and supercommutator
matrix rows open.  Hence the proposition proves the closure theorem for
supplied finite bialgebra data; the present compact \(K3\times E\)
source still has no populated primitive Hall bracket.

\begin{proposition}[The primitive Hall supercommutator satisfies
graded Jacobi; \textit{proved}]
\label{prop:finite-hall-graded-jacobi}
Let \(A\) be an associative \(\mathbb Z/2\)-graded algebra and let
\[
[x,y]=xy-(-1)^{|x||y|}yx
\]
for homogeneous \(x,y\in A\).  Then the bracket is graded
skew-symmetric and satisfies the graded Jacobi identity
\[
(-1)^{|x||z|}[x,[y,z]]
+(-1)^{|y||x|}[y,[z,x]]
+(-1)^{|z||y|}[z,[x,y]]=0 .
\]
Consequently, if the finite Hall bialgebra data of
Proposition~\ref{prop:hall-bialgebra-defects-primitive-closure} are
supplied so that \(P_R\) is closed under the supercommutator, then
\(P_R\) is a finite graded Lie algebra object under the Hall bracket.
\end{proposition}

\begin{proof}
Graded skew-symmetry is immediate from the definition.  For Jacobi,
expand the three displayed summands in the associative algebra \(A\).
The twelve monomials \(xyz,xzy,yxz,yzx,zxy,zyx\), each with its
Koszul sign, occur in cancelling pairs after using only associativity
and the parity rule \(|xy|=|x|+|y|\).  Thus the supercommutator on
any associative graded algebra is a graded Lie bracket.  If
\(x,y,z\in P_R\), Proposition~\ref{prop:hall-bialgebra-defects-primitive-closure}
keeps each inner and outer commutator in \(P_R\), so the same identity
holds in the primitive Hall subspace.
\end{proof}

The packet \(\texttt{certificates/hall/finite\_hall\_graded\_jacobi}\)
verifies
\(\texttt{FINITE\_HALL\_GRADED\_JACOBI\_OBSTRUCTION\_VERIFIED}\): it
imports the row \(354\) primitive-commutator closure packet, the row
\(350\) bialgebra-compatibility packet, and the compact Hall source
obstruction ledger; it records
Proposition~\ref{prop:finite-hall-graded-jacobi}; and it keeps the
source product matrix, primitive bracket matrix, associativity rows,
primitive-closure rows, Jacobi defect rows, parity rows, and
normal-ordered bracket matrices open.  Hence row \(355\) proves
graded Jacobi for supplied associative finite Hall data; it does not
populate a compact \(K3\times E\) primitive bracket.

\begin{proposition}[The Hall bracket preserves normal-ordered Gram
degree; \textit{proved for supplied homogeneous data}]
\label{prop:finite-hall-normal-ordered-gram-degree}
Let \(P_R=\bigoplus_{\widehat c\in\widehat\Gamma_R}P_{R,\widehat c}\)
be a finite primitive Hall stage whose bracket is obtained from a
homogeneous Hall product:
\[
[P_{R,\widehat c},P_{R,\widehat c'}]\subset
P_{R,\widehat c+\widehat c'} .
\]
Let
\[
\overline\Pi_X:\widehat\Gamma_R\longrightarrow\Gamma_R^\Pi
\]
be the additive normal-ordered Gram homomorphism.  Define the
descended fibre-summed pieces
\[
P^\Pi_{R,\gamma}=
\bigoplus_{\overline\Pi_X(\widehat c)=\gamma}P_{R,\widehat c}.
\]
Then
\[
[P^\Pi_{R,\gamma},P^\Pi_{R,\gamma'}]\subset
P^\Pi_{R,\gamma+\gamma'} .
\]
Equivalently, every non-zero finite bracket row has zero
normal-ordered Gram-degree defect
\[
\overline\Pi_X(\widehat e)
-\overline\Pi_X(\widehat c)
-\overline\Pi_X(\widehat c')=0
\]
for its source degrees \(\widehat c,\widehat c'\) and target degree
\(\widehat e=\widehat c+\widehat c'\).
\end{proposition}

\begin{proof}
The Hall product is homogeneous by hypothesis, so both ordered products
\(xy\) and \(yx\) in the supercommutator have source degree
\(\widehat c+\widehat c'\).  Subtracting them with the Koszul sign
cannot change the charge degree.  Applying the additive
normal-ordered homomorphism gives
\[
\overline\Pi_X(\widehat c+\widehat c')
=\overline\Pi_X(\widehat c)+\overline\Pi_X(\widehat c').
\]
If several source charges have the same descended Gram label, the
definition of \(P^\Pi_{R,\gamma}\) is the direct sum over that fibre;
the same additivity sends the whole bracket of the two fibres into the
fibre over \(\gamma+\gamma'\).  This proves the displayed inclusion.
\end{proof}

The packet
\(\texttt{certificates/hall/finite\_hall\_normal\_ordered\_gram\_degree}\)
verifies
\(\texttt{FINITE\_HALL\_NORMAL\_ORDERED\_GRAM\_DEGREE\_OBSTRUCTION\_VERIFIED}\):
it imports the row \(355\) graded-Jacobi packet, the finite
\(\texttt{hall\_pairing\_pushforward\_compatibility}\) packet, the
empty \(K3\times E\) charge-window fixture, and the compact Hall source
obstruction ledger; it records
Proposition~\ref{prop:finite-hall-normal-ordered-gram-degree}; and it
keeps the current source degree rows, normal-ordered lift rows, Hall
product rows, primitive bracket rows, and bracket-degree defect rows
open.  Hence row \(356\) proves the Gram-degree statement for supplied
homogeneous normal-ordered Hall data; it does not populate the compact
\(K3\times E\) bracket or replace the row \(357\) parity check.

\begin{proposition}[The Hall bracket respects parity;
\textit{proved for supplied homogeneous data}]
\label{prop:finite-hall-bracket-parity}
Let \(A=A_{\bar0}\oplus A_{\bar1}\) be a
\(\mathbb Z/2\)-graded associative algebra whose product is even:
\[
A_pA_q\subset A_{p+q}\qquad(p,q\in\mathbb Z/2).
\]
For homogeneous \(x\in A_p\) and \(y\in A_q\), set
\[
[x,y]=xy-(-1)^{pq}yx .
\]
Then \([x,y]\in A_{p+q}\).  Consequently, if the finite primitive Hall
stage \(P_R\) is supplied as a parity-graded subspace closed under this
supercommutator, then
\[
[P_{R,\bar p},P_{R,\bar q}]\subset P_{R,\bar p+\bar q}.
\]
After normal-ordered Gram descent, the same statement holds fibrewise:
\[
[P^\Pi_{R,\gamma,\bar p},P^\Pi_{R,\gamma',\bar q}]
\subset P^\Pi_{R,\gamma+\gamma',\bar p+\bar q}
\]
whenever the row \(356\) degree hypotheses and the parity block rows
are supplied.
\end{proposition}

\begin{proof}
Since the multiplication is even, both \(xy\) and \(yx\) lie in the
same parity summand \(A_{p+q}\).  Their signed difference therefore
also lies in \(A_{p+q}\).  Restricting to \(P_R\) uses only the supplied
closure of \(P_R\) under the supercommutator and the supplied parity
decomposition of \(P_R\).  The normal-ordered statement is obtained by
combining this parity calculation with
Proposition~\ref{prop:finite-hall-normal-ordered-gram-degree}: the
Gram label adds independently of the parity label, while the parity
label adds in \(\mathbb Z/2\).
\end{proof}

The packet \(\texttt{certificates/hall/finite\_hall\_bracket\_parity}\)
verifies
\(\texttt{FINITE\_HALL\_BRACKET\_PARITY\_OBSTRUCTION\_VERIFIED}\): it
imports the row \(355\) graded-Jacobi packet, the row \(356\)
normal-ordered Gram-degree packet, the finite
\(\texttt{hall\_pairing\_pushforward\_compatibility}\) bracket rows,
the finite \(\texttt{parity\_pushforward}\) packet, and the compact
Hall source obstruction ledger; it records
Proposition~\ref{prop:finite-hall-bracket-parity}; and it keeps the
current source parity blocks, homogeneous product rows, primitive
bracket rows, and bracket-parity defect rows open.  Hence row \(357\)
proves parity compatibility for supplied homogeneous finite Hall data;
it does not prove full \(K3\times E\) source parity or the GN--Kac
parity comparison required later in primitive recognition.

\begin{definition}[Compact geometric Hall bialgebra datum]
\label{def:compact-geometric-hall-bialgebra-datum}
At a finite HN height \(R\), a compact geometric Hall bialgebra datum is
the following finite correspondence package on the cosection-reduced
d-critical heart:
\begin{enumerate}
\item finite-type retained substacks
\(\mathfrak M^{\mathrm{red}}_{R,\widehat c}\) for all
\(\widehat c\in\widehat\Gamma_R\), with finite residual inertia,
reduced vanishing-cycle sheaves \(\Phi^{\mathrm{red}}_{R,\widehat c}\),
and Joyce--Upmeier orientation lines \(o_{R,\widehat c}\);
\item compactified extension correspondences
\[
\mathfrak M^{\mathrm{red}}_{R,\widehat c}\times
\mathfrak M^{\mathrm{red}}_{R,\widehat c'}
\xleftarrow{\ p\ }
\overline{\mathfrak E}^{\mathrm{red}}_{R,\widehat c,\widehat c'}
\xrightarrow{\ q\ }
\mathfrak M^{\mathrm{red}}_{R,\widehat c+\widehat c'}
\]
with admissible \(p^*,q_!\) for the product and admissible
\(q^*,p_!\) for the collision coproduct, in either a retained
compactified flag-Quot model or an explicitly supplied compact-support
six-functor model;
\item reduced Thom--Sebastiani isomorphisms
\(\operatorname{TS}^{\mathrm{red}}_{\widehat c,\widehat c'}\) and their
inverse transports, compatible with the multiplicative orientation
cocycle;
\item vacuum correspondences for the zero object, giving the unit
\(\eta_R\) and counit \(\epsilon_R\);
\item three-step and four-step retained flag stacks carrying the
associativity, coassociativity, and bialgebra-compatibility diagrams,
with base-change, projection-formula, Koszul sign, and
Thom--Sebastiani coherence witnesses;
\item reduced \(E\)-quotient and HN-transition compatibilities for all
the preceding correspondences.
\end{enumerate}
The product and coproduct attached to this datum are
\[
m_R=q_!\operatorname{TS}^{\mathrm{red}}p^*,\qquad
\Delta_R=p_!(\operatorname{TS}^{\mathrm{red}})^{-1}q^*.
\]
\end{definition}

\begin{definition}[Hall-product-preserving transition datum]
\label{def:transition-hall-product-preservation}
Let \(R'\le R\).  Suppose the finite-moduli transition geometry of
Definition~\ref{def:finite-moduli-transition-geometry}, the
orientation-preserving transition datum of
Definition~\ref{def:transition-orientation-preservation}, and the
vanishing-cycle-preserving transition datum of
Definition~\ref{def:transition-vanishing-cycle-preservation} are
supplied.  A Hall-product-preserving transition datum consists of:
\begin{enumerate}
\item a transition cube for the retained compactified extension
correspondences
\[
\mathfrak M_R\times\mathfrak M_R
\xleftarrow{\ p_R\ }\overline{\mathfrak E}_R
\xrightarrow{\ q_R\ }\mathfrak M_R,
\qquad
\mathfrak M_{R'}\times\mathfrak M_{R'}
\xleftarrow{\ p_{R'}\ }\overline{\mathfrak E}_{R'}
\xrightarrow{\ q_{R'}\ }\mathfrak M_{R'},
\]
whose input and output squares commute with the stage transitions;
\item transport of the coefficient system
\(\Phi_R^{\mathrm{red}}\otimes o_R\) on the extension graph, including
the Thom--Sebastiani product isomorphism
\[
\operatorname{TS}^{\mathrm{red}}_R:
p_R^*(\Phi_R^{\mathrm{red}}\boxtimes\Phi_R^{\mathrm{red}})
\xrightarrow{\sim}q_R^*\Phi_R^{\mathrm{red}};
\]
\item proper base-change, projection-formula, and compact-support
pushforward isomorphisms comparing the two pull-push functors
\[
q_{R,!}\operatorname{TS}^{\mathrm{red}}_Rp_R^*
\quad\text{and}\quad
q_{R',!}\operatorname{TS}^{\mathrm{red}}_{R'}p_{R'}^*
\]
after applying the transition maps;
\item finite product matrices \(M_R,M_{R'}\) in compact source bases
and transition matrices \(T_R,T_{R'}\) satisfying the product
intertwining identity
\[
T_{\alpha+\beta}\,M_R^{\alpha,\beta}
=M_{R'}^{\alpha,\beta}(T_\alpha\otimes T_\beta)
\]
in every retained degree and parity block;
\item compatibility with the associativity square on the retained
two-step flag stack;
\item zero correspondence, coefficient-system, base-change,
projection-formula, product-intertwining, associativity-transition, and
composition defect ranks, together with \(R^1\varprojlim=0\) for the
Hall product maps.
\end{enumerate}
Scalar traces, denominator products, signed multiplicities, target BKM
products, Pfaffian products, orientation transport alone, and
vanishing-cycle transport alone are not Hall-product transition data.
\end{definition}

\begin{proposition}[Transition preservation of compact Hall products;
\textit{proved}]
\label{prop:transition-hall-product-preservation-criterion}
Given a Hall-product-preserving transition datum of
Definition~\ref{def:transition-hall-product-preservation}, the
finite-stage transition \(R\to R'\) preserves the compact Hall product:
on every retained finite degree block,
\[
T_R\circ m_R=m_{R'}\circ(T_R\otimes T_R).
\]
\end{proposition}

\begin{proof}
The transition cube identifies the input pullbacks \(p_R^*\) and
\(p_{R'}^*\) after passage to the transition graph.  The coefficient
transport identifies \(\Phi^{\mathrm{red}}\otimes o\) and the
Thom--Sebastiani product isomorphisms on the two sides.  Proper base
change, the projection formula, and compact-support pushforward
compatibility then identify the two exceptional pushforwards along
\(q_R\) and \(q_{R'}\).  Hence the two pull-push functors defining
\(m_R\) and \(m_{R'}\) commute with transition.  In finite compact
source bases this is exactly the displayed matrix identity.  The
associativity-transition and \(R^1\varprojlim\) rows give strict
composition and inverse-limit exactness for the product maps.
\end{proof}

The packet
\(\texttt{certificates/hall/transition\_hall\_product\_preservation}\)
records the present state of this datum.  Its verifier status is
\texttt{TRANSITION\_HALL\_PRODUCT\_PRESERVATION\_OBSTRUCTION\_VERIFIED}:
the transition-geometry input, orientation-transition input,
vanishing-cycle transition input, source product matrices, extension
correspondence transition squares, input and output product-stack
squares, coefficient-system transport rows, Thom--Sebastiani product
transport rows, proper base-change rows, projection-formula rows,
compact-support pushforward rows, product-matrix intertwining rows,
associativity-transition rows, and Hall-product Mittag--Leffler rows
are not yet supplied.  Hence transition preservation of compact Hall
products remains open.

\begin{definition}[Hall-coproduct-preserving transition datum]
\label{def:transition-hall-coproduct-preservation}
Let \(R'\le R\).  Suppose the finite-moduli transition geometry of
Definition~\ref{def:finite-moduli-transition-geometry}, the
orientation-preserving transition datum of
Definition~\ref{def:transition-orientation-preservation}, and the
vanishing-cycle-preserving transition datum of
Definition~\ref{def:transition-vanishing-cycle-preservation} are
supplied.  A Hall-coproduct-preserving transition datum consists of:
\begin{enumerate}
\item a transition cube for the same retained compactified extension
correspondences, read in the splitting direction
\[
\mathfrak M_R
\xleftarrow{\ q_R\ }\overline{\mathfrak E}_R
\xrightarrow{\ p_R\ }\mathfrak M_R\times\mathfrak M_R,
\qquad
\mathfrak M_{R'}
\xleftarrow{\ q_{R'}\ }\overline{\mathfrak E}_{R'}
\xrightarrow{\ p_{R'}\ }\mathfrak M_{R'}\times\mathfrak M_{R'},
\]
whose \(q\)-source square and \(p\)-two-output square commute with the
stage transitions;
\item transport of the coefficient system
\(\Phi_R^{\mathrm{red}}\otimes o_R\) on the splitting graph, including
the inverse Thom--Sebastiani coproduct isomorphism
\[
(\operatorname{TS}^{\mathrm{red}}_R)^{-1}:
q_R^*\Phi_R^{\mathrm{red}}
\xrightarrow{\sim}
p_R^*(\Phi_R^{\mathrm{red}}\boxtimes\Phi_R^{\mathrm{red}});
\]
\item proper base-change, projection-formula, and compact-support
pushforward isomorphisms comparing the two pull-push functors
\[
p_{R,!}(\operatorname{TS}^{\mathrm{red}}_R)^{-1}q_R^*
\quad\text{and}\quad
p_{R',!}(\operatorname{TS}^{\mathrm{red}}_{R'})^{-1}q_{R'}^*
\]
after applying the transition maps;
\item finite coproduct matrices \(D_R,D_{R'}\) in compact source bases,
counit matrices \(\epsilon_R,\epsilon_{R'}\), and transition matrices
\(T_\gamma\) satisfying the coproduct intertwining identity
\[
(T_\mu\otimes T_\nu)\,D_R^{\rho;\mu,\nu}
=D_{R'}^{\rho;\mu,\nu}\,T_\rho
\]
in every retained degree and parity block with \(\rho=\mu+\nu\);
\item compatibility with the coassociativity square on the retained
two-step flag stack and with the zero-object counit square;
\item zero correspondence, coefficient-system, base-change,
projection-formula, coproduct-intertwining, coassociativity-transition,
counit-transition, and composition defect ranks, together with
\(R^1\varprojlim=0\) for the Hall coproduct maps.
\end{enumerate}
Scalar traces, denominator products, signed multiplicities, target BKM
coalgebra data, Pfaffian operations, bar-coalgebra deconcatenation,
Hall-product preservation, orientation transport alone, and
vanishing-cycle transport alone are not Hall-coproduct transition data.
\end{definition}

\begin{proposition}[Transition preservation of compact Hall coproducts;
\textit{proved}]
\label{prop:transition-hall-coproduct-preservation-criterion}
Given a Hall-coproduct-preserving transition datum of
Definition~\ref{def:transition-hall-coproduct-preservation}, the
finite-stage transition \(R\to R'\) preserves the compact Hall
coproduct: on every retained finite degree block with \(\rho=\mu+\nu\),
\[
(T_\mu\otimes T_\nu)\circ\Delta_R^{\rho;\mu,\nu}
=\Delta_{R'}^{\rho;\mu,\nu}\circ T_\rho .
\]
\end{proposition}

\begin{proof}
The transition cube identifies the source pullbacks \(q_R^*\) and
\(q_{R'}^*\) after passage to the transition graph.  The coefficient
transport identifies \(\Phi^{\mathrm{red}}\otimes o\) and the inverse
Thom--Sebastiani isomorphisms on the two sides.  Proper base change,
the projection formula, and compact-support pushforward compatibility
then identify the two exceptional pushforwards along \(p_R\) and
\(p_{R'}\).  Hence the two pull-push functors defining \(\Delta_R\) and
\(\Delta_{R'}\) commute with transition.  In finite compact source
bases this is exactly the displayed matrix identity.  The
coassociativity-transition, counit-transition, and \(R^1\varprojlim\)
rows give strict composition, counital compatibility, and inverse-limit
exactness for the coproduct maps.
\end{proof}

The packet
\(\texttt{certificates/hall/transition\_hall\_coproduct\_preservation}\)
records the present state of this datum.  Its verifier status is
\texttt{TRANSITION\_HALL\_COPRODUCT\_PRESERVATION\_OBSTRUCTION\_VERIFIED}:
the transition-geometry input, orientation-transition input,
vanishing-cycle transition input, source coproduct matrices, source
counit maps, splitting correspondence transition squares, source object
and target product stack squares, coefficient-system transport rows,
inverse Thom--Sebastiani coproduct transport rows, proper base-change
rows, projection-formula rows, compact-support pushforward rows,
coproduct-matrix intertwining rows, coassociativity-transition rows,
counit-transition rows, and Hall-coproduct Mittag--Leffler rows are not
yet supplied.  Hence transition preservation of compact Hall coproducts
remains open.

\begin{definition}[Primitive-subspace-preserving transition datum]
\label{def:transition-primitive-subspace-preservation}
Let \(R'\le R\), and suppose the finite Hall stages carry counital
coproducts \((\Delta_R,\epsilon_R)\) and
\((\Delta_{R'},\epsilon_{R'})\).  Put
\[
\widetilde\Delta_R
=\Delta_R-\eta_R\epsilon_R\otimes\mathrm{id}
-\mathrm{id}\otimes\eta_R\epsilon_R,
\qquad
P_R=\ker\widetilde\Delta_R\cap\ker\epsilon_R .
\]
A primitive-subspace-preserving transition datum consists of:
\begin{enumerate}
\item finite compact-source matrices for
\(\Delta_R,\Delta_{R'},\epsilon_R,\epsilon_{R'}\), the associated
reduced coproducts \(\widetilde\Delta_R,\widetilde\Delta_{R'}\), and
the ambient Hall transition \(T_R:\mathcal H_R\to\mathcal H_{R'}\);
\item source and target primitive kernel rows, including inclusions
\(\iota_R:P_R\hookrightarrow\mathcal H_R\),
\(\iota_{R'}:P_{R'}\hookrightarrow\mathcal H_{R'}\) and finite
primitive projection matrices;
\item the reduced-coproduct and counit intertwining identities
\[
\widetilde\Delta_{R'}T_R=(T_R\otimes T_R)\widetilde\Delta_R,\qquad
\epsilon_{R'}T_R=\epsilon_R;
\]
\item a restricted transition matrix
\[
T_R^{\Prim}:P_R\longrightarrow P_{R'}
\]
satisfying \(\iota_{R'}T_R^{\Prim}=T_R\iota_R\), together with zero
primitive-image, projection-commutator, and composition defect ranks.
\end{enumerate}
Hall product preservation, scalar trace data, denominator products,
target root spaces, bar-coalgebra primitives, normal-ordered degree
maps, and primitive ranks without inclusions are not
primitive-subspace transition data.
\end{definition}

\begin{proposition}[Transition preservation of primitive Hall
subspaces; \textit{proved}]
\label{prop:transition-primitive-subspace-preservation-criterion}
Given a primitive-subspace-preserving transition datum of
Definition~\ref{def:transition-primitive-subspace-preservation}, the
ambient Hall transition carries primitives to primitives:
\[
T_R(P_R)\subseteq P_{R'}.
\]
Equivalently, \(T_R|_{P_R}\) factors through the displayed
\(T_R^{\Prim}:P_R\to P_{R'}\).
\end{proposition}

\begin{proof}
Let \(x\in P_R\).  Then
\(\widetilde\Delta_Rx=0\) and \(\epsilon_Rx=0\).  The two
intertwining identities give
\[
\widetilde\Delta_{R'}(T_Rx)
=(T_R\otimes T_R)\widetilde\Delta_Rx=0,\qquad
\epsilon_{R'}(T_Rx)=\epsilon_Rx=0.
\]
Thus \(T_Rx\in\ker\widetilde\Delta_{R'}\cap\ker\epsilon_{R'}=P_{R'}\).
The finite restricted matrix \(T_R^{\Prim}\) is the same statement in
the chosen primitive bases.  The projection and composition defect
rows assert that this factorization is compatible with the finite
primitive summands and with successive stage transitions.
\end{proof}

The packet
\(\texttt{certificates/hall/transition\_primitive\_subspace\_preservation}\)
records the present state of this datum.  Its verifier status is
\texttt{TRANSITION\_PRIMITIVE\_SUBSPACE\_PRESERVATION\_OBSTRUCTION\_VERIFIED}:
the source and target coproduct matrices, source and target counit
maps, reduced-coproduct matrices, primitive kernel rows, ambient Hall
transition matrices, Hall-coproduct transition input, counit-transition
compatibility, kernel-intertwining rows, primitive projection matrices,
restricted primitive transition matrices, and primitive-transition
composition rows are not yet supplied.  Hence transition preservation
of primitive Hall subspaces remains open.

\begin{proposition}[Geometric Hall data give the finite bialgebra
defect rows; \textit{proved relative to the datum}]
\label{prop:geometric-hall-data-give-bialgebra-defects}
Let the compact geometric Hall bialgebra datum of
Definition~\ref{def:compact-geometric-hall-bialgebra-datum} be given at
height \(R\).  Then the matrices of
\((m_R,\Delta_R,\eta_R,\epsilon_R)\) satisfy the unit, counit,
associativity, coassociativity, and bialgebra-compatibility equations
of Definition~\ref{def:finite-hall-bialgebra-defect-tuple}.  Hence the
first seven components of
\(\mathfrak O_{\mathrm{Hbialg},R}\) vanish, and
Proposition~\ref{prop:hall-bialgebra-defects-primitive-closure}
forms the primitive Hall bracket.  Conversely, finite matrices
\((M,D,\eta,\epsilon)\) whose bialgebra defect ranks vanish but which
are not obtained from such a compact geometric Hall bialgebra datum are
not compact \(K3\times E\) source rows; they are only an abstract finite
bialgebra.
\end{proposition}

\begin{proof}
The left and right unit identities are the pull-push identities for the
vacuum extension correspondences.  The counit identities are their
collision analogues after applying the augmentation to the zero-object
factor.  Associativity of \(m_R\) is the equality of the two pull-push
composites on the retained three-step extension flag stack; the
base-change isomorphism, projection formula, and triple
Thom--Sebastiani coherence identify the two composites.  The same
argument applied to the reversed correspondence gives coassociativity
of \(\Delta_R\).

Bialgebra compatibility is the comparison between first multiplying and
then splitting an extension, and first splitting the two inputs and then
multiplying the four resulting pieces.  The retained four-step
flag-of-flags stack is the common refinement of these two operations.
Base change on this common refinement gives the equality of the two
pull-push functors; the middle interchange contributes exactly the
Koszul sign braiding \(\tau\), and the multiplicative
Thom--Sebastiani orientation transport identifies the coefficient
systems.  This is the matrix identity
\[
\Delta_Rm_R
=(m_R\otimes m_R)(1\otimes\tau\otimes1)
(\Delta_R\otimes\Delta_R).
\]
Thus all first seven defect matrices have rank zero.  The primitive
closure is then the algebraic calculation of
Proposition~\ref{prop:hall-bialgebra-defects-primitive-closure}.

The converse sentence is the Beilinson cut: an abstract finite
bialgebra with the same matrices is not a \(K3\times E\) source object
unless its entries are the pull-push matrices of the retained compact
correspondences with the stated orientation, quotient, and transition
data.
\end{proof}

\begin{lemma}[Primitive bracket lifting; \textit{conditional on the
Frobenius defect $\mathfrak o_F=0$}]
\label{lem:primitive-bracket-lifting}
Let $\langle\cdot,\cdot\rangle_X$ be the protected Hopf pairing on
$\widetilde P_X$, non-degenerate on the radical quotient.  The
trilinear Frobenius defect
\[
\mathfrak o_F(x,y,z):=\langle[x,y],z\rangle
+(-1)^{|x|\cdot|y|}\langle y,[x,z]\rangle,
\quad
x,y,z\in\widetilde P_X,\quad c_x+c_y+c_z=0,
\]
viewed as a class in
$C^3_{\mathrm{cyc}}(\widetilde P_X;k)$ after the reduced $E$-quotient
and radical quotient, vanishes precisely when the bracket is
graded-symmetric for the Hopf pairing.  Vanishing is the cyclic-CY-3
Serre-functor identity of
Theorem~\ref{thm:hall-product}.  Non-degeneracy of the
Hopf pairing alone does \emph{not} force $\mathfrak o_F=0$; it only
detects the vanishing once supplied.
\end{lemma}

\begin{definition}[Finite positive-negative Hall pairing]
\label{def:finite-hall-positive-negative-pairing}
Let
\[
P_R^\Pi=\bigoplus_{\gamma\in\Gamma_R^\Pi,\ \bar p\in\mathbb Z/2}
P_{R,\gamma,\bar p}^\Pi
\]
be a finite normal-ordered primitive Hall stage, equipped with the
descended product \(m_R^\Theta\) and a degree-zero trace
\[
\operatorname{tr}_{0,R}:
H_{R,0,\bar0}^{\Pi}\longrightarrow k .
\]
A finite positive-negative Hall pairing is the collection of bilinear
forms
\[
\langle\cdot,\cdot\rangle_{R,\gamma,\bar p}:
P_{R,\gamma,\bar p}^\Pi\otimes
P_{R,-\gamma,\bar p}^\Pi\longrightarrow k
\]
defined, in homogeneous finite bases, by the matrix
\[
G_{\gamma,\bar p;ab}
=\operatorname{tr}_{0,R}
\bigl(m_R^\Theta(e_{\gamma,\bar p,a},e_{-\gamma,\bar p,b})\bigr).
\]
Equivalently, \(G_{\gamma,\bar p;ab}\) is the pull-push coefficient of
the compact positive-negative pairing correspondence
\[
\mathfrak P^{\mathrm{red}}_{R,\gamma,-\gamma;0}
\longrightarrow
\mathfrak M^{\mathrm{red}}_{R,0}
\]
with reduced Thom--Sebastiani and orientation transport, followed by
\(\operatorname{tr}_{0,R}\).  A finite row of type \(G\) consists of
the degree \(\gamma\), the parity \(\bar p\), the positive and negative
primitive basis labels, the scalar coefficient, the pairing
correspondence, and the geometric source reference.

This is a definition of the source matrix \(G\).  It does not assert
that \(G\) is homogeneous, graded-symmetric, invariant under the Hall
bracket, compatible with the coproduct, or non-degenerate after the
radical quotient.  Those are the separate defect rows of
Definition~\ref{def:finite-hopf-pairing-identity-tuple}.
\end{definition}

The packet
\(\texttt{certificates/hall/finite\_hall\_positive\_negative\_pairing}\)
verifies
\(\texttt{FINITE\_HALL\_POSITIVE\_NEGATIVE\_PAIRING\_OBSTRUCTION\_VERIFIED}\):
it imports the formal degree-zero pairing rows of
\(\texttt{hall\_pairing\_pushforward\_compatibility}\), the row \(356\)
normal-ordered Gram-degree packet, the row \(357\) parity packet, and
the compact Hall source obstruction ledger; it records
Definition~\ref{def:finite-hall-positive-negative-pairing}; and it
keeps the current source \(G\)-entries, degree-zero trace row,
positive-negative pairing correspondences, orientation transport rows,
and Hopf-pairing identity rows open.  Hence row \(358\) defines the
finite positive-negative Hall pairing matrix; it does not prove any of
the row \(359\)--\(362\) pairing identities.

\begin{proposition}[The finite Hall pairing is homogeneous;
\textit{proved for supplied source data}]
\label{prop:finite-hall-pairing-homogeneity}
Assume the finite positive-negative Hall pairing of
Definition~\ref{def:finite-hall-positive-negative-pairing} is supplied
from a homogeneous normal-ordered Hall product and a degree-zero trace.
Extend it by zero to all pairs
\[
P_{R,\gamma,\bar p}^{\Pi}\otimes P_{R,\gamma',\bar q}^{\Pi}.
\]
Then
\[
\langle
P_{R,\gamma,\bar p}^{\Pi},
P_{R,\gamma',\bar q}^{\Pi}
\rangle_R=0
\quad\text{unless}\quad
\gamma+\gamma'=0\ \text{and}\ \bar p=\bar q .
\]
Equivalently, every non-zero finite pairing row has total
normal-ordered Gram degree \(0\) and belongs to one parity block.
\end{proposition}

\begin{proof}
For homogeneous inputs \(x\in P_{R,\gamma,\bar p}^{\Pi}\) and
\(y\in P_{R,\gamma',\bar q}^{\Pi}\), the row \(356\) degree
calculation gives
\[
m_R^\Theta(x,y)\in H_{R,\gamma+\gamma',\bar p+\bar q}^{\Pi}.
\]
The trace in Definition~\ref{def:finite-hall-positive-negative-pairing}
is defined on the degree-zero block.  Hence
\(\operatorname{tr}_{0,R}(m_R^\Theta(x,y))\) can contribute only when
\(\gamma+\gamma'=0\).  The positive-negative pairing matrix \(G\) is
defined blockwise on
\(P_{R,\gamma,\bar p}^{\Pi}\otimes P_{R,-\gamma,\bar p}^{\Pi}\); its
extension by zero kills the off-parity blocks.  Thus every non-zero
entry has total normal-ordered Gram degree \(0\) and one fixed parity
label.
\end{proof}

The packet
\(\texttt{certificates/hall/finite\_hall\_pairing\_homogeneity}\)
verifies
\(\texttt{FINITE\_HALL\_PAIRING\_HOMOGENEITY\_OBSTRUCTION\_VERIFIED}\):
it imports the row \(358\) positive-negative pairing definition, the
row \(356\) normal-ordered Gram-degree packet, the formal degree-zero
pairing rows of
\(\texttt{hall\_pairing\_pushforward\_compatibility}\), and the compact
Hall source obstruction ledger; it records
Proposition~\ref{prop:finite-hall-pairing-homogeneity}; and it keeps
the current source \(G\)-entries, trace row, pairing correspondence
rows, and homogeneity defect rows open.  Hence row \(359\) proves the
homogeneity criterion for supplied finite source data; it does not
prove graded supersymmetry, Frobenius invariance, coproduct
adjointness, or quotient non-degeneracy.

\begin{proposition}[The finite Hall pairing is supersymmetric;
\textit{proved for supplied signed-transpose rows}]
\label{prop:finite-hall-pairing-supersymmetry}
Assume the finite positive-negative Hall pairing of
Definition~\ref{def:finite-hall-positive-negative-pairing} is
homogeneous in the sense of
Proposition~\ref{prop:finite-hall-pairing-homogeneity}.  Suppose that
for every retained degree \(\gamma\), parity \(\bar p\), and pair of
finite basis labels \(a,b\), the signed-transpose row
\[
G_{\gamma,\bar p;ab}
-(-1)^{\bar p}G_{-\gamma,\bar p;ba}=0
\]
is supplied.  Then the extended finite pairing is supersymmetric:
\[
\langle x,y\rangle_R
=(-1)^{|x||y|}\langle y,x\rangle_R
\]
for all homogeneous \(x,y\in P_R^\Pi\).

Conversely, a compact-source proof of this identity requires the
source \(G\)-entries on both signs, source parity blocks,
signed-transpose defect rows, and the reduced Calabi--Yau
threefold Serre-sign and orientation-transport data.  Formal
degree-zero pairing rows, target positive-negative blocks, and scalar
traces do not supply these signs.
\end{proposition}

\begin{proof}
Write \(x\in P_{R,\gamma,\bar p}^{\Pi}\) and
\(y\in P_{R,\gamma',\bar q}^{\Pi}\).  By
Proposition~\ref{prop:finite-hall-pairing-homogeneity}, both
\(\langle x,y\rangle_R\) and \(\langle y,x\rangle_R\) vanish unless
\(\gamma'=-\gamma\) and \(\bar q=\bar p\).  On a non-zero block one
therefore has \(|x|=|y|=\bar p\), so
\((-1)^{|x||y|}=(-1)^{\bar p}\).  Expanding in the chosen finite
bases, the asserted supersymmetry equation is exactly
\[
G_{\gamma,\bar p;ab}
=(-1)^{\bar p}G_{-\gamma,\bar p;ba}.
\]
Linearity gives the identity for arbitrary homogeneous vectors.  This
is a finite signed-transpose criterion; the geometric computation of
the signs is the source Serre/orientation problem named in the final
sentence of the statement.
\end{proof}

The packet
\(\texttt{certificates/hall/finite\_hall\_pairing\_supersymmetry}\)
verifies
\(\texttt{FINITE\_HALL\_PAIRING\_SUPERSYMMETRY\_OBSTRUCTION\_VERIFIED}\):
it imports the row \(359\) homogeneity packet, the row \(358\)
positive-negative pairing definition, the reduced Calabi--Yau
threefold Serre-sign criterion packet, and the compact Hall source
obstruction ledger; it records
Proposition~\ref{prop:finite-hall-pairing-supersymmetry}; and it keeps
the current source \(G\)-entries, source parity blocks,
signed-transpose rows, Serre-sign rows, and orientation-transport rows
open.  Hence row \(360\) proves the finite supersymmetry criterion for
supplied source data; it does not prove Frobenius invariance, coproduct
adjointness, or quotient non-degeneracy.

\begin{definition}[Finite Hopf-pairing identity tuple]
\label{def:finite-hopf-pairing-identity-tuple}
Fix homogeneous finite bases on a retained primitive stage and write
\((M,B,D,G,K,Q)\) for the product, bracket, coproduct, Hopf-pairing,
radical, and quotient-splitting matrices.  The finite Hopf-pairing
identity tuple is
\[
\mathfrak O_{\mathrm{Hpair},R}
=\bigl(
\mathfrak o^{\mathrm{adj}}_R,\,
\mathfrak o^{\mathrm{Frob}}_R,\,
\mathfrak o^{\mathrm{qnd}}_R
\bigr),
\]
where
\[
\mathfrak o^{\mathrm{adj}}_R(x,y,z)
=\langle m_R^\Theta(x,y),z\rangle_R
-\langle x\otimes y,\Delta_R^\Theta z\rangle_R,
\]
\[
\mathfrak o^{\mathrm{Frob}}_R(x,y,z)
=\langle[x,y]_\Theta,z\rangle_R
+(-1)^{|x||y|}\langle y,[x,z]_\Theta\rangle_R,
\]
and \(\mathfrak o^{\mathrm{qnd}}_R\) is the finite rank defect of the
pairing induced by \(G\) on
\((P_R^\Pi/\operatorname{Rad}_R^\Pi)_\gamma\otimes
(P_R^\Pi/\operatorname{Rad}_R^\Pi)_{-\gamma}\) in every retained
parity block.  In matrix form these are the adjointness, Frobenius
cyclic, and quotient non-degeneracy rows of the finite source proof
tables.
\end{definition}

\begin{proposition}[The finite Hall pairing is invariant;
\textit{proved for supplied Frobenius rows}]
\label{prop:finite-hall-pairing-invariance}
Assume the finite primitive Hall bracket is the graded-skew
supercommutator and is parity-compatible, the finite Hall pairing is
supersymmetric in the sense of
Proposition~\ref{prop:finite-hall-pairing-supersymmetry}, and the
Frobenius cyclic component of
Definition~\ref{def:finite-hopf-pairing-identity-tuple} vanishes:
\[
\mathfrak o^{\mathrm{Frob}}_R(x,y,z)=0
\]
for every homogeneous retained triple.  Then the finite pairing is
invariant under the Hall bracket:
\[
\langle [x,y]_\Theta,z\rangle_R
=\langle x,[y,z]_\Theta\rangle_R
\]
for all homogeneous \(x,y,z\in P_R^\Pi\).

Conversely, a compact-source proof of this identity requires source
primitive bracket rows, source \(G\)-entries, parity blocks, Frobenius
defect rows, cyclic three-point pairing correspondences, and the
reduced Calabi--Yau threefold Serre-sign/orientation transport data.
The graded Jacobi identity, pairing supersymmetry, or target invariant
form does not supply these source rows.
\end{proposition}

\begin{proof}
Put \(p=|x|\), \(q=|y|\), \(r=|z|\).  Apply the assumed vanishing of
\(\mathfrak o^{\mathrm{Frob}}_R\) to the triple \((y,z,x)\):
\[
0=\langle [y,z]_\Theta,x\rangle_R
  +(-1)^{qr}\langle z,[y,x]_\Theta\rangle_R .
\]
The bracket parity row gives
\(|[y,z]_\Theta|=q+r\) and \(|[x,y]_\Theta|=p+q\).  Supersymmetry of
the pairing gives
\[
\langle [y,z]_\Theta,x\rangle_R
=(-1)^{p(q+r)}\langle x,[y,z]_\Theta\rangle_R .
\]
Using graded skew-symmetry
\([y,x]_\Theta=-(-1)^{pq}[x,y]_\Theta\) and supersymmetry again,
\[
\langle z,[y,x]_\Theta\rangle_R
=-(-1)^{pq+(p+q)r}
\langle [x,y]_\Theta,z\rangle_R .
\]
The two exponents \(p(q+r)\) and \(qr+pq+(p+q)r\) agree in
\(\mathbb Z/2\).  After removing the common sign, the displayed
vanishing becomes
\[
\langle x,[y,z]_\Theta\rangle_R
-\langle [x,y]_\Theta,z\rangle_R=0,
\]
which is the asserted invariance.
\end{proof}

The packet
\(\texttt{certificates/hall/finite\_hall\_pairing\_invariance}\)
verifies
\(\texttt{FINITE\_HALL\_PAIRING\_INVARIANCE\_OBSTRUCTION\_VERIFIED}\):
it imports the row \(360\) supersymmetry packet, the row \(357\)
bracket-parity packet, the row \(355\) graded-Jacobi packet, the
reduced Calabi--Yau threefold Serre-sign criterion packet, and the
compact Hall source obstruction ledger; it records
Proposition~\ref{prop:finite-hall-pairing-invariance}; and it keeps
the current source \(B\)-entries, \(G\)-entries, parity blocks,
Frobenius defect rows, cyclic three-point correspondences,
Serre-sign rows, and orientation-transport rows open.  Hence row
\(361\) proves the finite invariance criterion for supplied source
data; it does not prove coproduct adjointness or quotient
non-degeneracy.

\begin{proposition}[The finite Hall product and coproduct are adjoint;
\textit{proved for supplied Hopf-adjointness rows}]
\label{prop:finite-hall-pairing-coproduct-adjointness}
Assume the finite Hall product \(m_R^\Theta\), coproduct
\(\Delta_R^\Theta\), and positive-negative pairing \(G\) are supplied
in homogeneous finite bases.  Equip \(P_R^\Pi\otimes P_R^\Pi\) with the
Koszul tensor product pairing
\[
\langle x_1\otimes x_2,u_1\otimes u_2\rangle_{R\otimes R}
=(-1)^{|x_2||u_1|}
\langle x_1,u_1\rangle_R\langle x_2,u_2\rangle_R .
\]
If every adjointness component of
Definition~\ref{def:finite-hopf-pairing-identity-tuple} vanishes,
\[
\mathfrak o^{\mathrm{adj}}_R(x,y,z)=0 ,
\]
then
\[
\langle m_R^\Theta(x,y),z\rangle_R
=\langle x\otimes y,\Delta_R^\Theta z\rangle_{R\otimes R}
\]
for all homogeneous retained finite vectors \(x,y,z\) in the domain of
the Hopf-pairing tuple.

In coordinates, for homogeneous basis vectors
\(e_a,e_b,e_d\), this is the rank-zero matrix identity
\[
\sum_c M_{ab}^{c}G_{cd}
-\sum_{u,v}D_{d}^{uv}(-1)^{|e_b||e_u|}G_{au}G_{bv}=0 .
\]
A compact-source proof therefore requires source \(M\)-entries,
\(D\)-entries, \(G\)-entries, tensor-pairing sign rows, Hopf
adjointness defect rows, and the product/coproduct correspondence
orientation data.  Bialgebra compatibility, pairing invariance, target
Hopf forms, and scalar traces do not supply this adjointness matrix.
\end{proposition}

\begin{proof}
Write
\[
m_R^\Theta(e_a,e_b)=\sum_c M_{ab}^{c}e_c,\qquad
\Delta_R^\Theta(e_d)=\sum_{u,v}D_d^{uv}e_u\otimes e_v .
\]
By bilinearity of \(G\),
\[
\langle m_R^\Theta(e_a,e_b),e_d\rangle_R
=\sum_c M_{ab}^{c}G_{cd}.
\]
By the Koszul tensor product rule,
\[
\langle e_a\otimes e_b,\Delta_R^\Theta(e_d)\rangle_{R\otimes R}
=\sum_{u,v}D_d^{uv}(-1)^{|e_b||e_u|}G_{au}G_{bv}.
\]
Thus the displayed coordinate defect is exactly
\(\mathfrak o^{\mathrm{adj}}_R(e_a,e_b,e_d)\).  Its vanishing on the
finite homogeneous basis gives the identity for basis vectors, and
linearity gives it for all homogeneous vectors.
\end{proof}

The packet
\(\texttt{certificates/hall/finite\_hall\_pairing\_coproduct\_adjointness}\)
verifies
\(\texttt{FINITE\_HALL\_PAIRING\_COPRODUCT\_ADJOINTNESS\_OBSTRUCTION\_VERIFIED}\):
it imports the row \(361\) pairing-invariance packet, the row \(342\)
product-operator packet, the row \(347\) coproduct-operator packet, the
row \(350\) bialgebra-compatibility packet, and the compact Hall source
obstruction ledger; it records
Proposition~\ref{prop:finite-hall-pairing-coproduct-adjointness}; and
it keeps the current source \(M\)-entries, \(D\)-entries,
\(G\)-entries, tensor-pairing sign rows, Hopf-adjointness defect rows,
product/coproduct correspondence rows, and orientation-transport rows
open.  Hence row \(362\) proves the coproduct-adjointness criterion for
supplied finite source data; it does not prove quotient
non-degeneracy.

\begin{proposition}[The finite quotient Hopf pairing is non-degenerate;
\textit{proved for supplied quotient-rank rows}]
\label{prop:finite-hall-quotient-pairing-nondegeneracy}
Let \(G_{\gamma,\bar p}\) be a supplied homogeneous finite
positive-negative pairing matrix on
\[
P_{R,\gamma,\bar p}^{\Pi}\otimes P_{R,-\gamma,\bar p}^{\Pi}.
\]
Let \(K_{\gamma,\bar p}\subset P_{R,\gamma,\bar p}^{\Pi}\) and
\(K_{-\gamma,\bar p}\subset P_{R,-\gamma,\bar p}^{\Pi}\) be supplied
left and right radical bases for \(G_{\gamma,\bar p}\), and let
\[
Q_{\gamma,\bar p}:L_{\gamma,\bar p}\longrightarrow
P_{R,\gamma,\bar p}^{\Pi},\qquad
Q_{-\gamma,\bar p}:L_{-\gamma,\bar p}\longrightarrow
P_{R,-\gamma,\bar p}^{\Pi}
\]
be quotient splittings.  Write
\[
\overline G_{\gamma,\bar p}
=Q_{\gamma,\bar p}^{t}G_{\gamma,\bar p}Q_{-\gamma,\bar p}
\]
for the induced quotient pairing matrix.  If the finite rank rows give
\[
\operatorname{rank}\overline G_{\gamma,\bar p}
=\dim L_{\gamma,\bar p}
=\dim L_{-\gamma,\bar p}
\]
for every retained \((\gamma,\bar p)\), then the induced pairing
\[
L_{\gamma,\bar p}\otimes L_{-\gamma,\bar p}\longrightarrow k
\]
is non-degenerate.  This is the quotient non-degeneracy component
\(\mathfrak o^{\mathrm{qnd}}_R=0\) of
Definition~\ref{def:finite-hopf-pairing-identity-tuple}.

Conversely, a compact-source proof requires the source \(G\)-entries,
left and right radical bases \(K_{\pm\gamma,\bar p}\), quotient
splittings \(Q_{\pm\gamma,\bar p}\), quotient-pairing matrices, and
rank or determinant rows proving full rank.  Hopf adjointness,
Frobenius invariance, target radical tables, rank equality alone, and
signed dimensions do not supply quotient non-degeneracy.
\end{proposition}

\begin{proof}
Because \(K_{\gamma,\bar p}\) and \(K_{-\gamma,\bar p}\) are supplied
as the left and right radicals of \(G_{\gamma,\bar p}\), the pairing
descends to
\[
P_{R,\gamma,\bar p}^{\Pi}/K_{\gamma,\bar p}
\quad\text{and}\quad
P_{R,-\gamma,\bar p}^{\Pi}/K_{-\gamma,\bar p}.
\]
The splittings \(Q_{\gamma,\bar p}\) and \(Q_{-\gamma,\bar p}\) give
finite bases of these quotients, and the matrix of the descended
pairing in these bases is precisely
\(\overline G_{\gamma,\bar p}
=Q_{\gamma,\bar p}^{t}G_{\gamma,\bar p}Q_{-\gamma,\bar p}\).  Full
rank equal to both quotient dimensions says that the two adjoint maps
\[
L_{\gamma,\bar p}\longrightarrow L_{-\gamma,\bar p}^{*},
\qquad
L_{-\gamma,\bar p}\longrightarrow L_{\gamma,\bar p}^{*}
\]
are isomorphisms.  This is non-degeneracy of the quotient pairing.
\end{proof}

The packet
\(\texttt{certificates/hall/finite\_hall\_quotient\_pairing\_nondegeneracy}\)
verifies
\(\texttt{FINITE\_HALL\_QUOTIENT\_PAIRING\_NONDEGENERACY\_OBSTRUCTION\_VERIFIED}\):
it imports the row \(362\) coproduct-adjointness packet, the row \(360\)
supersymmetry packet, the row \(358\) positive-negative pairing
definition, the pairing-kernel \(\varprojlim^1\) obstruction packet,
the transition-radical obstruction packet, and the compact Hall source
obstruction ledger; it records
Proposition~\ref{prop:finite-hall-quotient-pairing-nondegeneracy}; and
it keeps the current source \(G\)-entries, \(K\)-entries,
\(Q\)-entries, quotient-pairing matrices, quotient-rank rows,
determinant rows, and radical-identification rows open.  Hence row
\(363\) proves the quotient non-degeneracy criterion for supplied
finite source data; it does not prove transition preservation of
radicals or \(\varprojlim^1\)-vanishing for the radical tower.

\begin{definition}[Compact geometric Hopf-pairing datum]
\label{def:compact-geometric-hopf-pairing-datum}
At finite HN height \(R\), assume the compact geometric Hall
bialgebra datum of
Definition~\ref{def:compact-geometric-hall-bialgebra-datum}.  A
compact geometric Hopf-pairing datum on the descended primitive stage is
the following additional source data:
\begin{enumerate}
\item a degree-zero protected trace
\[
\operatorname{tr}_{0,R}:H_c^\bullet
(\mathfrak M^{\mathrm{red}}_{R,0},\Phi^{\mathrm{red}}_{R,0}\otimes
o_{R,0})\longrightarrow k
\]
compatible with the vacuum counit and with HN transitions;
\item positive-negative compact pairing correspondences
\(\mathfrak P^{\mathrm{red}}_{R,\gamma,-\gamma;0}\) obtained by
restricting the Hall product to total descended Gram degree \(0\), with
orientation and Thom--Sebastiani transport inherited from the Hall
datum;
\item homogeneous pairing matrices
\[
G_{\gamma,ab}
=\operatorname{tr}_{0,R}\!\left(
m_R(e_{\gamma,a},e_{-\gamma,b})\right)
\]
on every retained parity block, together with kernel matrices
\(K_{\gamma,\bar p}\) and quotient splittings \(Q_{\gamma,\bar p}\);
\item triple pairing correspondences for charges
\(\gamma_1+\gamma_2+\gamma_3=0\), with cyclic permutation
isomorphisms induced by the \(3\)-Calabi--Yau Serre pairing and by
Joyce--Upmeier orientation transport;
\item finite rank-zero witnesses for Hopf adjointness and the primitive
Frobenius cyclic identity, and finite rank witnesses that the induced
pairing on
\((P_R^\Pi/\operatorname{Rad}_R^\Pi)_\gamma\otimes
(P_R^\Pi/\operatorname{Rad}_R^\Pi)_{-\gamma}\)
is non-degenerate in every retained parity block.
\end{enumerate}
The radical is
\[
\operatorname{Rad}_{R,\gamma,\bar p}^{\Pi}
=\{x\in P_{R,\gamma,\bar p}^{\Pi}\mid
\langle x,y\rangle_R=0\text{ for all }
y\in P_{R,-\gamma,\bar p}^{\Pi}\},
\]
with both signs included.  The datum is source-side: the target
Gritsenko--Nikulin radical may be used only as a comparison codomain
after the source radical and quotient have been formed.
\end{definition}

The packet
\(\texttt{certificates/hall/finite\_hall\_compact\_geometric\_hopf\_pairing\_datum}\)
verifies
\(\texttt{FINITE\_HALL\_COMPACT\_GEOMETRIC\_HOPF\_PAIRING\_DATUM\_OBSTRUCTION\_VERIFIED}\):
it imports the row \(363\) quotient non-degeneracy packet, the row
\(362\) coproduct-adjointness packet, the row \(361\) Frobenius
invariance packet, the reduced Calabi--Yau threefold Serre-sign
criterion packet, the level-\(\mathsf Z\) protected-trace obstruction
ledger as a firewall, and the compact Hall source obstruction ledger;
it records
Definition~\ref{def:compact-geometric-hopf-pairing-datum}; and it
keeps the source degree-zero trace, positive-negative pairing
correspondences, orientation and Thom--Sebastiani transport,
\(G\)-entries, \(K\)-entries, \(Q\)-entries, cyclic three-point
correspondences, Hopf-adjointness witnesses, Frobenius cyclic
witnesses, and quotient non-degeneracy witnesses open.  Hence row
\(364\) records the compact geometric Hopf-pairing datum as a
source-level datum; it does not construct it, and the
level-\(\mathsf Z\) protected trace does not substitute for the source
degree-zero trace.

\begin{definition}[Radical-preserving transition datum]
\label{def:transition-radical-preservation}
Let \(R'\le R\).  Suppose the primitive-subspace transition datum of
Definition~\ref{def:transition-primitive-subspace-preservation} is
given after normal-ordered descent, so that
\[
T_R^{\Prim}:P_R^\Pi\longrightarrow P_{R'}^\Pi
\]
is defined.  A radical-preserving transition datum consists of:
\begin{enumerate}
\item source and target Hopf-pairing matrices
\(G_{R,\gamma,\bar p}\) and \(G_{R',\gamma,\bar p}\), with the Hopf
adjointness, Frobenius cyclic, and quotient non-degeneracy rows of
Definition~\ref{def:finite-hopf-pairing-identity-tuple};
\item source and target radical kernel rows
\[
K_{R,\gamma,\bar p}=\operatorname{Rad}_{R,\gamma,\bar p}^{\Pi},
\qquad
K_{R',\gamma,\bar p}=\operatorname{Rad}_{R',\gamma,\bar p}^{\Pi},
\]
together with quotient splittings \(Q_{R,\gamma,\bar p}\) and
\(Q_{R',\gamma,\bar p}\);
\item for every target opposite primitive \(y'\in
P_{R',-\gamma,\bar p}^{\Pi}\), a lifted source opposite primitive
\(L_{R'R}^{-\gamma}y'\in P_{R,-\gamma,\bar p}^{\Pi}\) with
\(T_R^{\Prim,-\gamma}L_{R'R}^{-\gamma}y'=y'\);
\item the pairing-transition identity
\[
\bigl\langle T_R^{\Prim,\gamma}x,y'\bigr\rangle_{R'}
=\bigl\langle x,L_{R'R}^{-\gamma}y'\bigr\rangle_R
\]
for every \(x\in P_{R,\gamma,\bar p}^{\Pi}\) and
\(y'\in P_{R',-\gamma,\bar p}^{\Pi}\);
\item a restricted radical transition matrix
\[
T_R^{\operatorname{Rad}}:K_{R,\gamma,\bar p}\longrightarrow
K_{R',\gamma,\bar p}
\]
and a descended quotient transition matrix
\[
\bar T_R:P_{R,\gamma,\bar p}^{\Pi}/K_{R,\gamma,\bar p}
\longrightarrow
P_{R',\gamma,\bar p}^{\Pi}/K_{R',\gamma,\bar p},
\]
with zero radical-image, quotient well-definedness, splitting
compatibility, and composition defect ranks.
\end{enumerate}
Scalar traces, denominator products, signed multiplicities, target
radical tables, pairing matrices without Hopf identities, rank
equalities, primitive transition alone, and normal-ordered degree maps
are not radical-transition data.
\end{definition}

\begin{proposition}[Transition preservation of Hopf radicals;
\textit{proved}]
\label{prop:transition-radical-preservation-criterion}
Given a radical-preserving transition datum of
Definition~\ref{def:transition-radical-preservation}, the primitive
transition carries the Hopf-pairing radical into the Hopf-pairing
radical:
\[
T_R^{\Prim,\gamma}\bigl(\operatorname{Rad}_{R,\gamma,\bar p}^{\Pi}\bigr)
\subseteq
\operatorname{Rad}_{R',\gamma,\bar p}^{\Pi}.
\]
Consequently \(T_R^{\Prim}\) descends to the radical quotients through
the displayed \(\bar T_R\).
\end{proposition}

\begin{proof}
Let \(x\in\operatorname{Rad}_{R,\gamma,\bar p}^{\Pi}\), and let
\(y'\in P_{R',-\gamma,\bar p}^{\Pi}\).  Choose the supplied lift
\(L_{R'R}^{-\gamma}y'\in P_{R,-\gamma,\bar p}^{\Pi}\).  The
pairing-transition identity gives
\[
\bigl\langle T_R^{\Prim,\gamma}x,y'\bigr\rangle_{R'}
=\bigl\langle x,L_{R'R}^{-\gamma}y'\bigr\rangle_R=0,
\]
because \(x\) pairs trivially with every source opposite primitive.
Since this holds for every \(y'\), the vector
\(T_R^{\Prim,\gamma}x\) lies in
\(\operatorname{Rad}_{R',\gamma,\bar p}^{\Pi}\).  The finite restricted
matrix \(T_R^{\operatorname{Rad}}\) records this inclusion in radical bases, and the
quotient well-definedness row is precisely the statement that
\(T_R^{\Prim}\) factors through \(\bar T_R\).  The splitting and
composition rows make the quotient transition independent of the
chosen finite representatives and compatible with successive stages.
\end{proof}

The packet
\(\texttt{certificates/hall/transition\_radical\_preservation}\)
records the present state of this datum.  Its verifier status is
\texttt{TRANSITION\_RADICAL\_PRESERVATION\_OBSTRUCTION\_VERIFIED}: the
primitive-transition input, source and target Hopf-pairing matrices,
source and target Hopf-pairing identity rows, source and target radical
kernels, opposite-primitive lift rows, pairing-transition identities,
radical-image identities, restricted radical transition matrices,
source and target quotient splittings, descended quotient transition
matrices, and radical-transition composition rows are not yet supplied.
Hence transition preservation of Hopf radicals remains open.

\begin{proposition}[Geometric Hopf-pairing data give the finite
Hopf-pairing defect rows; \textit{proved relative to the datum}]
\label{prop:geometric-hopf-pairing-data-give-defects}
Let the compact geometric Hall bialgebra datum and the compact
geometric Hopf-pairing datum be given at height \(R\).  Then
\[
\mathfrak O_{\mathrm{Hpair},R}=0.
\]
Consequently the radical of the source pairing is a coproduct coideal;
if the cyclic Frobenius component of the datum is imposed on primitives,
it is also a Lie ideal, and the quotient carries the descended primitive
Lie bracket and Hopf pairing.  Conversely, a matrix \(G\), a kernel
matrix \(K\), and a quotient matrix \(Q\) with the right ranks are not a
compact source Hopf-pairing datum unless their entries are obtained
from the protected trace, pairing correspondences, and cyclic
three-point coherences above.
\end{proposition}

\begin{proof}
Hopf adjointness is the equality of two pull-push evaluations on the
same three-point extension/collision correspondence.  On one side one
multiplies \(x\) and \(y\), then pairs with \(z\) by
\(\operatorname{tr}_{0,R}\); on the other one splits \(z\) by
\(\Delta_R\), pairs the two resulting factors with \(x\) and \(y\), and
then applies the tensor trace.  The common flag refinement and the
projection formula identify the two Borel--Moore chain maps, while the
Thom--Sebastiani and orientation transports identify the coefficient
systems.  This is
\[
\langle m_R^\Theta(x,y),z\rangle_R
=\langle x\otimes y,\Delta_R^\Theta z\rangle_R.
\]

The Frobenius cyclic identity is the same three-point trace after a
cyclic permutation of the inputs.  The \(3\)-Calabi--Yau Serre pairing
changes the order by the Koszul sign; the supercommutator therefore
gives
\[
\langle[x,y]_\Theta,z\rangle_R
+(-1)^{|x||y|}\langle y,[x,z]_\Theta\rangle_R=0.
\]
The quotient non-degeneracy component is not a consequence of the first
two equations; it is the finite rank condition in the datum after
dividing by \(K_{\gamma,\bar p}\).  Hence the three components of
\(\mathfrak O_{\mathrm{Hpair},R}\) vanish.

The coideal and Lie-ideal conclusions are
Lemma~\ref{lem:hopf-radical-ideal-coideal}.  The converse sentence is a
source/target firewall: a finite bilinear form can have the same kernel
rank as the target radical without being the protected trace pairing of
the compact \(K3\times E\) Hall source.
\end{proof}

The packet
\(\texttt{certificates/hall/finite\_hall\_hopf\_pairing\_defect\_rows}\)
verifies
\(\texttt{FINITE\_HALL\_HOPF\_PAIRING\_DEFECT\_ROWS\_OBSTRUCTION\_VERIFIED}\):
it imports the row \(364\) compact geometric Hopf-pairing datum packet,
the row \(350\) Hall bialgebra compatibility packet, the row \(362\)
Hopf-adjointness criterion, the row \(361\) Frobenius cyclic criterion,
the row \(363\) quotient non-degeneracy criterion, and the compact Hall
source obstruction ledger; it records
Proposition~\ref{prop:geometric-hopf-pairing-data-give-defects}; and it
keeps the present compact source Hopf-adjointness rows, Frobenius
cyclic rows, quotient non-degeneracy rows, \(M\)-, \(D\)-, \(G\)-,
\(K\)-, and \(Q\)-entries, radical ideal/coideal rows, degree-zero
source trace, positive-negative pairing correspondences, and cyclic
three-point coherences open.  Hence row \(365\) proves only the
relative implication from supplied compact geometric Hopf-pairing data
to \(\mathfrak O_{\mathrm{Hpair},R}=0\); it does not construct the
datum and does not prove Hopf-radical descent for the current compact
source.

\begin{proposition}[A pairing matrix is not a Hopf-pairing proof;
\textit{proved}]
\label{prop:pairing-matrix-not-hopf-pairing-proof}
At a finite HN stage, the matrix \(G\), its kernel matrix \(K\), and a
quotient splitting \(Q\) do not by themselves prove that the radical is
the Hopf radical used in primitive recognition.  The finite stage proves
Hopf-radical descent only after
\(\mathfrak O_{\mathrm{Hpair},R}=0\): Hopf adjointness, the primitive
Frobenius cyclic identity, and non-degeneracy of the induced quotient
pairing.  Under these three vanishings, the radical is a coproduct
coideal and a Lie ideal, hence the quotient carries the descended
primitive Lie bracket and Hopf pairing.
\end{proposition}

\begin{proof}
The matrix \(G\) defines a bilinear form and \(K=\ker G\) records its
linear radical.  Coideal stability is not a linear-algebra consequence
of being a kernel: it follows from Hopf adjointness
\(\langle m(x,y),r\rangle=\langle x\otimes y,\Delta r\rangle\) and
non-degeneracy of the quotient tensor pairing.  Lie-ideal stability is
not a consequence of kernel membership either; it follows from the
Frobenius identity
\(\langle[x,r],z\rangle=-(-1)^{|x||r|}\langle r,[x,z]\rangle\).  These
are exactly the first two defects in
\(\mathfrak O_{\mathrm{Hpair},R}\), and quotient non-degeneracy is the
third.  With all three vanishing, Lemma~\ref{lem:hopf-radical-ideal-coideal}
applies.  If any one is absent, \(K\) is only the radical of a finite
matrix, not the Hopf radical required by
Theorem~\ref{thm:primitive-recognition}.
\end{proof}

The packet
\(\texttt{certificates/hall/finite\_hall\_pairing\_matrix\_firewall}\)
verifies
\(\texttt{FINITE\_HALL\_PAIRING\_MATRIX\_FIREWALL\_OBSTRUCTION\_VERIFIED}\):
it imports the row \(365\) Hopf-pairing defect packet, the row \(362\)
Hopf-adjointness criterion, the row \(361\) Frobenius cyclic criterion,
the row \(363\) quotient non-degeneracy criterion, and the compact Hall
source obstruction ledger; it records
Proposition~\ref{prop:pairing-matrix-not-hopf-pairing-proof}; and it
keeps the current source \(G\)-, \(K\)-, and \(Q\)-entries,
Hopf-adjointness rows, Frobenius cyclic rows, quotient
non-degeneracy rows, radical coideal rows, radical Lie-ideal rows, and
quotient Hopf-pairing rows open.  Hence row \(366\) is a firewall
theorem: \(G,K,Q\) and their ranks can identify a finite bilinear
kernel, but they do not prove the Hopf radical used in primitive
recognition.  Hopf-radical descent begins only after
\(\mathfrak O_{\mathrm{Hpair},R}=0\).

\begin{lemma}[Hopf radical is a Lie ideal and a coideal;
\textit{conditional on $\mathfrak o_F=0$ and Hopf adjointness}]
\label{lem:hopf-radical-ideal-coideal}
Let $P_R^\Pi$ be a finite normal-ordered primitive stage with
transported product $m_R^\Theta$, coproduct $\Delta_R^\Theta$, graded
symmetric Hopf pairing $\langle\cdot,\cdot\rangle_R$, and radical
$\operatorname{Rad}_R^\Pi$.  Assume Hopf adjointness
$\langle m_R^\Theta(x,y),z\rangle_R
=\langle x\otimes y,\Delta_R^\Theta z\rangle_R$
and non-degeneracy on
$P_R^\Pi/\operatorname{Rad}_R^\Pi\otimes
P_R^\Pi/\operatorname{Rad}_R^\Pi$.  Then
\[
(\bar q_{\Pi,R}\otimes\bar q_{\Pi,R})\Delta_R^\Theta(r)=0
\quad(r\in\operatorname{Rad}_R^\Pi),
\]
i.e.\ $\operatorname{Rad}_R^\Pi$ is a coideal.  If, in addition,
$\mathfrak o_{F,R}=0$, then $\operatorname{Rad}_R^\Pi$ is a Lie
ideal.
\end{lemma}

\begin{proof}
For $r\in\operatorname{Rad}_R^\Pi$ and $x,y\in P_R^\Pi$, Hopf
adjointness gives
$\langle x\otimes y,\Delta_R^\Theta(r)\rangle_R
=\langle m_R^\Theta(x,y),r\rangle_R=0$;
non-degeneracy of the quotient tensor pairing then kills the class.
The Lie-ideal half is direct: if $\mathfrak o_{F,R}=0$ then
$\langle[x,r]_\Theta,z\rangle_R
=-(-1)^{|x||r|}\langle r,[x,z]_\Theta\rangle_R=0$ for every $z$,
so $[x,r]_\Theta\in\operatorname{Rad}_R^\Pi$.
\end{proof}

The packet
\(\texttt{certificates/hall/finite\_hall\_hopf\_radical\_ideal\_coideal}\)
verifies
\(\texttt{FINITE\_HALL\_HOPF\_RADICAL\_IDEAL\_COIDEAL\_OBSTRUCTION\_VERIFIED}\):
it imports the row \(366\) pairing-matrix firewall, the row \(365\)
Hopf-pairing defect packet, the row \(362\) Hopf-adjointness criterion,
the row \(361\) Frobenius cyclic criterion, the row \(363\) quotient
non-degeneracy criterion, and the compact Hall source obstruction
ledger; it records Lemma~\ref{lem:hopf-radical-ideal-coideal}; and it
keeps the present source Hopf-adjointness rows, Frobenius cyclic rows,
quotient tensor non-degeneracy rows, radical coideal rows, radical
Lie-ideal rows, \(G\)-, \(K\)-, and \(Q\)-entries open.  Hence row
\(367\) proves the conditional algebraic criterion for a supplied
finite Hopf pairing; it does not prove the radical ideal/coideal rows
for the current compact \(K3\times E\) source.

\subsubsection*{Integral form and stable-envelope transport}

Two parallel constructions act on the same compact Hall heart:
\begin{enumerate}
\item the integral-form comparison modeled on the
Schiffmann--Vasserot elliptic Hall algebra; a compact
\(K3\times E\) specialization would require a non-toric
Maulik--Okounkov-type stable-envelope transport;
\item the proposed Drinfeld double
\(\mathcal D_\hbar(\textup{CoHA}_{K3\times E})\) of Vol III, whose
positive half must be the protected primitive Hopf algebra above and
whose Drinfeld pairing must be the Hopf pairing.
\end{enumerate}
The compatibility datum between these two constructions is a comparison
with Calabi--Yau quantum groups.  A different constant or
sign would be a distinct mathematical claim; the Drinfeld double is not
asserted here as an input theorem.

\begin{definition}[Integral-form stable-envelope comparison datum]
\label{def:integral-form-stable-envelope-comparison-datum}
At finite height \(R\), an integral-form stable-envelope comparison
datum for the compact \(K3\times E\) Hall heart consists of the
following finite rows:
\begin{enumerate}
\item an integral lattice
\((Y^+_{R,\mathbb Z}(X),m_R,\Delta_R,\langle-,-\rangle_R)\)
inside the protected Hall bialgebra, stable under product,
coproduct, and the Hopf pairing;
\item a separated and complete topology on \(Y^+_R(X)\), compatible
with the HN transition maps and the Hopf radical quotient;
\item a completed Drinfeld double
\(\mathcal D_\hbar(Y^+_R(X))\) with cross-relations whose structure
constants are the finite Hall pairing rows;
\item a Cartan part in the double, identified with
\(\Lambda_{II}^{2,1}\otimes\mathbb C\), and rows proving that this
Cartan survives the reduced primitive quotient;
\item a compact non-toric stable-envelope correspondence satisfying
support, normalization, degree, and chamber-cocycle identities on the
retained \(K3\times E\) moduli;
\item comparison constants, signs, and degree shifts identifying the
preceding double with the Calabi--Yau quantum-group normalization.
\end{enumerate}
The datum is source-level.  The target BKM denominator, the
Schiffmann--Vasserot affine \(A_0\) theorem, and the
Maulik--Okounkov quiver-variety stable envelope do not supply these
finite compact rows.
\end{definition}

\begin{proposition}[Stable-envelope and Drinfeld-double scope
firewall; \textit{proved}]
\label{prop:integral-form-stable-envelope-transport-firewall}
The Maulik--Okounkov stable-envelope theorem and the
Schiffmann--Vasserot elliptic-Hall comparison may be used here only in
their stated scope.  They do not construct the integral form,
separated completion, compact non-toric stable envelope, Drinfeld
double cross-relations, or Cartan identification required for the
compact \(K3\times E\) Hall heart.  Any assertion
\[
\mathcal D_\hbar(Y^+(K3\times E))^{\mathrm{int}}_{\mathrm{compl}}
 \simeq
Y_\hbar(\widehat{\mathfrak g}_{\Delta_5})
\]
is therefore conditional on the datum of
Definition~\ref{def:integral-form-stable-envelope-comparison-datum}.
\end{proposition}

\begin{proof}
Appendix~\ref{app:stable-envelope} records the precise imported
theorems.  Maulik--Okounkov stable envelopes are constructed for
smooth symplectic varieties in the quiver-variety setting with a torus
action and chamber data; their theorem gives the geometric
\(R\)-matrix in that category.  Schiffmann--Vasserot identify the
\(\mathbb C^3\) cohomological Hall algebra with the positive half in
the affine \(A_0\) elliptic-Hall setting.  Neither theorem contains a
compact non-toric \(K3\times E\) moduli atlas, a retained D0-HN
comparison, an \(O2\) wall atlas, or the Hopf-radical quotient rows
above.

Thus the cited theorems fix the target shape and the normalization
problem, but they do not supply the finite source rows of
Definition~\ref{def:integral-form-stable-envelope-comparison-datum}.
The compact comparison is a separate datum: without those rows the
symbol \(\mathcal D_\hbar(Y^+(K3\times E))^{\mathrm{int}}_{\mathrm{compl}}\)
is a requested construction, not an established input theorem.
\end{proof}

The packet
\(\texttt{certificates/hall/finite\_hall\_integral\_form\_stable\_envelope\_transport}\)
verifies
\(\texttt{FINITE\_HALL\_INTEGRAL\_FORM\_STABLE\_ENVELOPE\_TRANSPORT\_OBSTRUCTION\_VERIFIED}\):
it imports the row \(367\) Hopf-radical ideal/coideal criterion, the
compact Hall source ledger, the D0-HN obstruction ledger, and the
\((O2)\) wall-atlas ledger; it records
Definition~\ref{def:integral-form-stable-envelope-comparison-datum}
and
Proposition~\ref{prop:integral-form-stable-envelope-transport-firewall};
and it keeps the present compact source integral-form rows, product
stability rows, coproduct stability rows, pairing stability rows,
separated-completion rows, complete-completion rows, Drinfeld-double
rows, cross-relation rows, stable-envelope rows, Cartan
identification rows, primitive Cartan survival rows, compact
non-toric replacement theorem, and Calabi--Yau quantum-group
comparison rows open.  Hence row \(368\) proves only the scope
firewall and the required finite datum; it does not construct the
compact \(K3\times E\) integral form, the compact Drinfeld double, or
stable-envelope transport.

\subsubsection*{Conilpotent source coalgebra and Hopf radical
quotient}

\begin{definition}[Finite conilpotent source-coalgebra datum]
\label{def:conilpotent-source-coalgebra}
A finite conilpotent source-coalgebra datum at HN height \(R\) is a
source-side package
\[
(C_{X,R},\eta^C_R,\varepsilon^C_R,
F^{\mathrm{co}}_{R,\bullet},
\Delta_R^{\mathrm{ch}},\rho^C_{R'R})
\]
obtained from the retained collision correspondences on
\(\mathcal F_{X,\sigma,S,\le R}^{\mathrm{hyb}}\), after the reduced
\(E\)-quotient, with the following finite rows:
\begin{enumerate}
\item \(C_{X,R}\) is a coaugmented chiral coalgebra on \(E\), with
coaugmentation \(\eta^C_R\), counit \(\varepsilon^C_R\), and
coalgebra identities for \(\Delta_R^{\mathrm{ch}}\);
\item \(C_{X,R}\) is filtered by a finite increasing source charge
filtration
\[
0=F^{\mathrm{co}}_{R,-1}\subset F^{\mathrm{co}}_{R,0}\subset
\cdots\subset F^{\mathrm{co}}_{R,N_R}=C_{X,R}
\]
whose associated graded pieces are finite dimensional in each
\(\Gamma_R^\Pi\)-degree, where
\(\Gamma_R^\Pi:=\overline\Pi_X(\widehat\Gamma_R)\)
is the image, under the \emph{normal-ordered} Gram homomorphism
\(\overline\Pi_X\) of Chapter~\ref{ch:k3xe-setup}, of the retained
normal-ordered charge window $\widehat\Gamma_R$ at Harder--Narasimhan
height \(R\) (\S\ref{subsec:hybrid-ran-carrier});
\item the additive \(\overline\Pi_X\), not the raw quadratic
\(\Pi_X\), is the grading map;
\item the reduced iterated coproduct is nilpotent:
\[
(\overline\Delta_R^{\mathrm{ch}})^{N_R+1}=0;
\]
\item for every successor \(R'\le R\), the transition map
\(\rho^C_{R'R}\) preserves the filtration, counit, coaugmentation, and
collision coproduct.
\end{enumerate}
Only after these transition rows and the strict Mittag--Leffler rows
are supplied does the inverse limit define the source coalgebra
\[
C_X=\varprojlim_R C_{X,R}.
\]
\end{definition}

\begin{proposition}[Conilpotent source-coalgebra firewall;
\textit{proved}]
\label{prop:conilpotent-source-coalgebra-firewall}
The notation \(C_{X,R}\), the formal bar-coalgebra construction, the
target BKM denominator, the stable-envelope comparison, and the
Drinfeld-double comparison do not construct the finite conilpotent
source-coalgebra datum of
Definition~\ref{def:conilpotent-source-coalgebra}.  A compact
\(K3\times E\) source coalgebra is available only after the
coaugmentation, counit, finite filtration, collision-coproduct,
coalgebra-identity, conilpotence, quotient-after-correspondence,
transition, and Mittag--Leffler rows are supplied.
\end{proposition}

\begin{proof}
A conilpotent coalgebra is not determined by its name.  It requires a
coalgebra object, counit and coaugmentation, a coproduct satisfying
coassociativity and counit identities, and a filtration for which the
reduced iterated coproduct terminates.  In this manuscript those
structures must be produced by retained compact collision
correspondences and their quotient-after-correspondence identities.
The target denominator and the later Drinfeld-double comparison live
on the comparison side; stable-envelope transport is a separate
normalization problem.  None of them supplies the source chain maps or
zero-defect rows listed in
Definition~\ref{def:conilpotent-source-coalgebra}.  The inverse-limit
object \(C_X\) additionally requires transition compatibility and
strict Mittag--Leffler exactness, since otherwise
\(R^1\varprojlim\) remains an obstruction.
\end{proof}

The packet
\(\texttt{certificates/hall/finite\_hall\_conilpotent\_source\_coalgebra}\)
verifies
\(\texttt{FINITE\_HALL\_CONILPOTENT\_SOURCE\_COALGEBRA\_OBSTRUCTION\_VERIFIED}\):
it imports the row \(368\) stable-envelope/Drinfeld-double firewall,
the hybrid carrier ledger, the compact Hall source obstruction
ledger, the D0-HN obstruction ledger, the Hall product and coproduct
transition ledgers, and the primitive-space and pairing-kernel
\(\varprojlim^1\) ledgers; it records
Definition~\ref{def:conilpotent-source-coalgebra} and
Proposition~\ref{prop:conilpotent-source-coalgebra-firewall}; and it
keeps the present compact source coalgebra rows, coaugmentation rows,
counit rows, finite filtration rows, collision-coproduct rows,
coalgebra-identity rows, conilpotence rows, quotient coalgebra
identity rows, transition coalgebra rows, inverse-limit
Mittag--Leffler rows, and bar-cobar comparison rows open.  Hence row
\(369\) proves only the source-data boundary for \(C_{X,R}\) and
\(C_X\); it does not construct the compact \(K3\times E\)
conilpotent source coalgebra.

\begin{proposition}[Sectorial Hall truncation and inverse limit;
\textit{conditional on
$[\textup{K3$\times$E-fin}]$, $[\textup{K3$\times$E-atlas}]$,
and $[\textup{ML}]$}]
\label{prop:sectorial-hall-truncation}
Under the three standing hypotheses
$[\textup{K3$\times$E-fin}]$ (finite-type semistable
substacks at bounded HN height),
$[\textup{K3$\times$E-atlas}]$ (finite d-critical/cosection
Hall atlas with Joyce--Upmeier orientation lines), and
$[\textup{ML}]$ (Mittag--Leffler descent for the HN inverse system),
the finite Hall product, coproduct, primitive projection, and Hopf
pairing are compatible with the transition maps $\rho_{R'R}$, and
their inverse limits define the completed Hall Hopf datum
$\widetilde P_X$, $C_X$.  The Mittag--Leffler hypothesis controls
the $R^1\varprojlim$ obstruction on every finite-window slice; without
it, only a compatible pro-object is obtained.
\end{proposition}

The hypothesis $[\textup{K3$\times$E-fin}]$ is not a
consequence of fibrewise K3 boundedness alone.  A precise sufficient
input is product-stability boundedness in the
Abramovich--Polishchuk/Liu sense
\cite{AbramovichPolishchukTStructures,LiuProductStability}: for a
rational K3 stability condition $\sigma_S=(\mathcal A_S,Z_S)$ and
Liu's product stability $\sigma_{s,t}$ on
$D^b(\operatorname{Coh}(S\times E))$, fixed-class boundedness
\[
\mathfrak M^{ss}_{\sigma_{s,t}}(\gamma)\text{ is finite type over }\C
\]
together with Liu's support property would imply
$[\textup{K3$\times$E-fin}]$ for bounded $h_S$-windows.
Piyaratne--Toda boundedness for threefold Bridgeland/BMT
\cite{PiyaratneTodaModuli3folds} does not automatically supply this
theorem; fibrewise Bayer--Macri boundedness alone does not either.

\subsubsection*{Wall-crossing and chamber dependence}

The lattice cocycle $B(c,c')$ of the dyonic Mukai construction is
chamber-independent.  The compact Hall side $D_0(X;\sigma,S)$ is
chamber-dependent: finite-type semistable substacks change under
wall-crossing, and the Kontsevich--Soibelman wall-crossing formula
\cite{KSCoHA} relates Hall counts in adjacent chambers.  If the
datum $D_0(X;\sigma,S)$ is supplied in chamber $\sigma$, then
$\overline\Pi_X$ remains a homomorphism and the recognition target
is $\mathfrak g_{\Delta_5}$ provided the chamber identifications hold
for the OP/DT scalar import; that the chosen $\sigma$ lies in the
OP/DT chamber, or that the wall-crossing transport is compatible
with $\overline\Pi_X$, is a separate conjectural input.  No
chamber-independence of the recognition target is asserted.

\subsubsection*{The Hall obstruction class}

The Hall lane carries the residuals
\[
\mathfrak O_{\mathrm{Hall},R}
=(\mathfrak o^{\mathrm{stage}}_R,\,
\mathfrak o^{\mathrm{conf}}_R,\,
\mathfrak o^{\mathrm{prop}}_R,\,
\mathfrak o^{\mathrm{TS}}_R,\,
\mathfrak O_{\mathrm{Hbialg},R},\,
\mathfrak O_{\mathrm{Hpair},R},\,
\mathfrak o^{\mathrm{ML}}_{R'R}).
\]
The first four entries assert the existence of the compact geometric
Hall bialgebra datum of
Definition~\ref{def:compact-geometric-hall-bialgebra-datum}: finite
retained substacks, compact-support functoriality, and coherent
Thom--Sebastiani orientation transport.  The tuple
\(\mathfrak O_{\mathrm{Hbialg},R}\) is then the algebraic defect tuple
of Definition~\ref{def:finite-hall-bialgebra-defect-tuple}; it is not
defined before product, coproduct, unit, and counit matrices have been
obtained from the compact correspondences.  The tuple
\(\mathfrak O_{\mathrm{Hpair},R}\) is the Hopf-pairing identity tuple of
Definition~\ref{def:finite-hopf-pairing-identity-tuple}; it is
independent of the bialgebra equations and is not defined before the
compact geometric Hopf-pairing datum of
Definition~\ref{def:compact-geometric-hopf-pairing-datum}: protected
trace, positive-negative pairing correspondences, cyclic three-point
coherences, source radical, quotient splitting, and quotient
non-degeneracy.  On the projection-finite stratum the wrapped product
is the classical d-critical CoHA product after cosection reduction; on
the wrapped stratum the construction is open and governed by the (D0)
obstruction.  PBW compatibility remains a separate recognition datum on
the descended primitive layer.

\subsubsection*{Comparison with Drinfeld doubling}

Vol III's Drinfeld doubling $Y^+(X)\to G(X)$ is the comparison target
for the compact source side; the integral form is
the chiral lift of the Schiffmann--Vasserot elliptic Hall algebra
after the stable-envelope and completion data are fixed.
Maulik--Okounkov stable-envelope transport supplies the integral
form on the projection-finite Ran stratum; the wrapped stratum is
the corresponding stratum-completion.  No black-box equality
$\mathrm{CoHA}=W_\infty$
is asserted: $\mathrm{CoHA}(\C^3)=Y^+(\mathfrak{gl}_1)$ in Vol III,
and $W_{1+\infty}$ appears only after Drinfeld doubling, centre,
vacuum evaluation.  The same discipline applies here: the protected
primitive Hopf algebra of $K3\times E$ becomes the BKM target only
after the Hopf radical quotient, the Frobenius identity, and the
no-extra-relations theorem of \S\ref{subsec:primitive-recognition}.

Under \(\mathfrak O_{\mathrm{Hall},R}=0\) and the Mittag--Leffler
transition datum, the Hall correspondences give a finite-stage
protected primitive Hopf datum on the cosection-twisted heart.  This
datum lives at the operator level and has no automatic comparison to
the scalar denominator
$\mathrm{den}(\autcor)=64^{-1}\Delta_5(2Z)$ imported in
Chapter~\ref{ch:view2-bkm}.  Closing the cross-level identity
\textcircled{2}$\leftrightarrow$\textcircled{3} requires a chain-level
scalar incarnation of the primitive structure: a Pfaffian line
$\mathscr L_{\mathrm{Pf},R}$ on the reduced moduli, a Dirac block
decomposition, and a first-order Pfaffian section whose
$\Sp_4(\Z)$-pullback recovers $\Delta_5$.  This is the fourth entry
$(\mathrm L 4)$, constructed in \S\ref{subsec:dirac-pfaffian}.
